%% Question bank chapter file - generated by bank_tools.py build (do not edit while the bank is still being built:
%% it is regenerated from the per-test data in Question Bank\data). Every question is one bankq block:
%%   \begin{bankq}[id=..., source={...}, paper=.., qno=.., type=..., subtopic={...}, answer={...}, answer_src=key|solved,
%%                 working={...}, difficulty=1-10, diff_parts={K S I C F}, diff_note={...},
%%                 group=<passage id or empty>, also_in={duplicate copies in other tests}]
%% difficulty = min(10, K+S+I+C+F-1): K concepts 1-3, S steps 1-4, I insight 0-2, C computation 0-2, F format 0-1
%%   ...question LaTeX, options, figures (\bfig / \bfigi -> img/)...
%%   \end{bankq}
%% Passages shared by several questions are bankpassage blocks; their questions carry group=<passage id>.
\documentclass[10pt]{article}
\usepackage[a4paper,margin=1.8cm]{geometry}
\usepackage[T1]{fontenc}
\usepackage{amsmath,amssymb,mathtools}
\usepackage{graphicx}
\usepackage[export]{adjustbox}
\usepackage{xcolor,needspace,keyval}
\usepackage{enumitem}
\usepackage{booktabs,array,tabularx}
\usepackage[hidelinks]{hyperref}
\setlength{\parindent}{0pt}
\setlength{\parskip}{3pt}
\makeatletter
\define@key{bq}{id}{\def\bq@id{#1}}\define@key{bq}{source}{\def\bq@src{#1}}\define@key{bq}{paper}{\def\bq@paper{#1}}
\define@key{bq}{qno}{\def\bq@qno{#1}}\define@key{bq}{type}{\def\bq@type{#1}}\define@key{bq}{subtopic}{\def\bq@sub{#1}}
\define@key{bq}{answer}{\def\bq@ans{#1}}\define@key{bq}{answer_src}{\def\bq@asrc{#1}}\define@key{bq}{working}{\def\bq@work{#1}}
\define@key{bq}{group}{\def\bq@group{#1}}\define@key{bq}{also_in}{\def\bq@also{#1}}
\define@key{bq}{difficulty}{\def\bq@diff{#1}}\define@key{bq}{diff_parts}{\def\bq@dparts{#1}}\define@key{bq}{diff_note}{\def\bq@dnote{#1}}
\newenvironment{bankq}[1][]{\def\bq@work{}\def\bq@also{}\def\bq@group{}\def\bq@diff{?}\def\bq@dparts{}\def\bq@dnote{}\setkeys{bq}{#1}%
  \par\needspace{4\baselineskip}\medskip\noindent{\small\textbf{\bq@id}\hfill\color{gray}[\bq@src, Paper \bq@paper, Q\bq@qno]
  \bq@type\ $\cdot$ \bq@sub\ $\cdot$ \textbf{D\bq@diff}}\par\nopagebreak\relax}
  {\par{\small\color{teal}Answer: \textbf{\bq@ans}\ifx\bq@work\@empty\else\ (solved: \bq@work)\fi
  \ifx\bq@also\@empty\else\par\color{gray}Also in: \bq@also\fi}\par}
\define@key{bp}{id}{\def\bp@id{#1}}\define@key{bp}{source}{\def\bp@src{#1}}\define@key{bp}{questions}{\def\bp@qs{#1}}
\newenvironment{bankpassage}[1][]{\setkeys{bp}{#1}\par\needspace{4\baselineskip}\bigskip\noindent
  {\small\textbf{Passage \bp@id}\hfill\color{gray}[\bp@src; questions \bp@qs]}\par\nopagebreak\relax\itshape}{\par}
\makeatother
\newcommand{\bfig}[2]{\par\smallskip{\centering\includegraphics[width=#1pt,max width=0.62\linewidth,max height=0.32\textheight]{img/#2}\par}\smallskip}
\newcommand{\bfigi}[2]{\includegraphics[width=#1pt,max width=\linewidth,max height=0.32\textheight,valign=c]{img/#2}}
\newcommand{\bmissing}[1]{\textit{[figure missing in the source conversion: \detokenize{#1}]}}
\newenvironment{qoptions}{\begin{itemize}[leftmargin=1.8em,itemsep=1pt,topsep=2pt,label={}]}{\end{itemize}}
\newcommand{\opt}[2]{\item[(#1)] #2}
\providecommand{\fwstart}{}\providecommand{\fwend}{}\providecommand{\qfullsplit}{}
\newcolumntype{L}[1]{>{\raggedright\arraybackslash}p{#1}}
\newcolumntype{C}[1]{>{\centering\arraybackslash}p{#1}}
\newcolumntype{R}[1]{>{\raggedleft\arraybackslash}p{#1}}
\providecommand{\sectionhdr}[1]{}\providecommand{\subsectionhdr}[1]{}


\begin{document}
{\Large\bfseries Ionic Equilibrium --- Question Bank}\par\smallskip
{\small Chemistry chapter C07. Questions from all converted test papers, grouped by source series. Answers: printed key where it matches the paper, otherwise solved (answer\_src=solved, with working). Exact repeats across papers appear once, with also\_in listing the other copies.}\par\medskip\hrule

\section*{Score-I papers}
\begin{bankq}[id=SI12-P1-CHEM-S1-Q06, source={Score-I - Test 12 (09 Mar 2022)}, paper=1, qno=6, type=multi, subtopic={Solubility product}, answer={C,D}, answer_src=key, difficulty=5, diff_parts={K2 S2 I0 C1 F1}, diff_note={concepts: solubility product, complex formation | steps: solubility in water; in KBr; in AgNO3; in NH3 using K = Ksp Kf and solve}]
Which is/are correct statement(s) about the solubility of AgCl(s).\\ {Given: $K_{sp(\text{AgCl})}=10^{-10}$; $K_f[\text{Ag(NH}_3)_2^+]=10^8$.}
\begin{qoptions}
\opt{A}{Solubility of AgCl in pure water is $10^{-5}$ gm/litre}
\opt{B}{Solubility of AgCl in 2M KBr is $10^{-5}$ gm/litre}
\opt{C}{Solubility of AgCl in 2M AgNO$_3$ is $5\times10^{-11}$ M}
\opt{D}{Solubility of AgCl in 2M NH$_3$ is 0.166 M}
\end{qoptions}
\end{bankq}

\section*{Score-II papers}
\begin{bankq}[id=SII02-P2-CHEM-S3-Q01, source={Score-II - Test 02 (28 Mar 2021)}, paper=2, qno=1, type=numerical, subtopic={Solubility product}, answer={0.085}, answer_src=key, difficulty=8, diff_parts={K3 S1 I2 C2 F1}, diff_note={concepts: standard problem in this chapter, solubility product, apply the relevant results | steps: apply the relevant results; compute}]
Silver ion forms a complex ion with thiosulphate ion, [Ag(S$_2$O$_3$)$_2$]$^{3-}$, with $K_f=2.8\times10^{13}$. How much AgBr ($K_{sp}=5.4\times10^{-13}$), in mole, will dissolve in 1 litre of 0.20 M Na$_2$S$_2$O$_3$ solution.

(Report your answer till two places after decimal) ($\sqrt{15.12}=3.88$)
\end{bankq}
\begin{bankq}[id=SII04-P1-CHEM-S1-Q05, source={Score-II - Test 04 (07 May 2021)}, paper=1, qno=5, type=single, subtopic={pH}, answer={C}, answer_src=key, difficulty=4, diff_parts={K2 S1 I1 C1 F0}, diff_note={concepts: standard problem in this chapter, ph | steps: apply the relevant results; compute}]
Some aqueous solutions are given below. Find the solution whose pH changes maximum by adding 10 millimoles amount of pure H$_2$SO$_4$. [$\log7=0.84$]

\begin{qoptions}
\opt{A}{100 mL, 0.1 M HCl $+$ 100 mL, 0.5 M Ca(OH)$_2$}
\opt{B}{100 mL, 0.1 M Ca(OH)$_2$ $+$ 100 mL, 0.4 M H$_2$S}
\opt{C}{100 mL, 0.1 M NaOH $+$ 100 mL, 0.2 M K$_2$SO$_4$}
\opt{D}{100 mL, 0.5 M CH$_3$COOH $+$ 100 mL, 0.1 M Ba(OH)$_2$}
\end{qoptions}
\end{bankq}
\begin{bankq}[id=SII04-P2-CHEM-S4-Q07, source={Score-II - Test 04 (07 May 2021)}, paper=2, qno=7, type=numerical, subtopic={Buffers}, answer={6}, answer_src=key, difficulty=8, diff_parts={K3 S1 I2 C2 F1}, diff_note={concepts: standard problem in this chapter, buffers, apply the relevant results | steps: apply the relevant results; compute}]
Out of following how many mixtures are buffers solution\\
(i) NaCN (0.5 mole) \hspace{1em} (ii) NaCN (0.5 moles) $+$ HCN (0.5 moles)\\
(iii) NH$_3$ (1 mole) $+$ NH$_4$Cl (1 mole) \hspace{1em} (iv) CH$_3$COOH (1 mole) $+$ NaOH (0.75 moles)\\
(v) NH$_4$Cl (1 mole) $+$ NaOH (0.75 moles) \hspace{1em} (vi) CH$_3$COONa (0.5 moles) $+$ HCl (0.25 mole)\\
(vii) H$_2$CO$_3$ (1 mole) $+$ NaOH (1.5 moles)
\end{bankq}
\begin{bankq}[id=SII06-P2-CHEM-S2-Q03, source={Score-II - Test 06 (05 Jul 2022)}, paper=2, qno=3, type=numerical, subtopic={Solubility product}, answer={5}, answer_src=key, difficulty=8, diff_parts={K3 S1 I2 C2 F1}, diff_note={concepts: standard problem in this chapter, solubility product, apply the relevant results | steps: apply the relevant results; compute}]
What limiting amount of Na$_2$CO$_3$ (in millimoles) is to be added to 150 ml of 0.1 M carbonated water to make its pH = 6.4? Given: $K_{a_1}$ and $K_{a_2}$ for H$_2$CO$_3$ = $4\times10^{-8}$ \& $4\times10^{-12}$ respectively.
\end{bankq}
\begin{bankq}[id=SII07-P1-CHEM-S1-Q19, source={Score-II - Test 07 (10 Jul 2022)}, paper=1, qno=19, type=single, subtopic={pH}, answer={C}, answer_src=key, difficulty=4, diff_parts={K2 S1 I1 C1 F0}, diff_note={concepts: standard problem in this chapter, ph | steps: apply the relevant results; compute}]
A solution is prepared by mixing 50.01 ml, 0.01 M, NaOH(aq.) and 50 ml, 0.005 M H$_2$SO$_4$(aq.). If few drops of phenolpthalein are now added into this solution, then find fraction (approx.) of this indicator present in its coloured form.\\
Given that :\\
Indicator constant of phenolpthalein is $2\times10^{-8}$\\
Ionic product of water is $10^{-14}$
\begin{qoptions}
\opt{A}{$\dfrac{1}{2}$}
\opt{B}{$\dfrac{1}{3}$}
\opt{C}{$\dfrac{2}{3}$}
\opt{D}{$\dfrac{10}{11}$}
\end{qoptions}
\end{bankq}
\begin{bankq}[id=SII08-P1-CHEM-S3-Q03, source={Score-II - Test 08 (05 Feb 2023)}, paper=1, qno=3, type=numerical, subtopic={Buffers}, answer={10.15}, answer_src=key, difficulty=8, diff_parts={K3 S1 I2 C2 F1}, diff_note={concepts: standard problem in this chapter, buffers, apply the relevant results | steps: apply the relevant results; compute}]
A solution is prepared by mixing 0.02 mol H$_2$CO$_3$, 0.01 mol NaHCO$_3$, 0.02 mol Na$_2$CO$_3$ and 0.02 mol NaOH in 100 mL of water. pH of the resulting solution is\_\_\_\_\_\_.\\
\relax[Given: $pK_{a_1}$ and $pK_{a_2}$ of H$_2$CO$_3$ are 6.37 and 10.32, respectively; $\log 2 = 0.30$, $\log 3 = 0.47$]
\end{bankq}
\begin{bankq}[id=SII08-P2-CHEM-S3-Q01, source={Score-II - Test 08 (05 Feb 2023)}, paper=2, qno=1, type=numerical, subtopic={pH}, answer={2}, answer_src=key, difficulty=8, diff_parts={K3 S1 I2 C2 F1}, diff_note={concepts: standard problem in this chapter, ph, apply the relevant results | steps: apply the relevant results; compute}]
Concentration of H$_2$SO$_4$ and Na$_2$SO$_4$ in a solution is 1 M and $2\times10^{-2}$ M, respectively. Molar solubility of PbSO$_4$ in the same solution is $X\times10^{-Y}$ M (expressed in scientific notation). The value of $X-Y$ is\_\_\_\_\_\_\_.\\
\relax[Given: Solubility product of PbSO$_4$ $(K_{sp}) = 1.6\times10^{-8}$. For H$_2$SO$_4$, $K_{a1}$ is very large and $K_{a2}=0.2\times10^{-2}$]
\end{bankq}
\begin{bankq}[id=SII09-P1-CHEM-S2-Q09, source={Score-II - Test 09 (12 Feb 2023)}, paper=1, qno=9, type=single, subtopic={pH}, answer={A}, answer_src=key, difficulty=4, diff_parts={K2 S1 I1 C1 F0}, diff_note={concepts: standard problem in this chapter, ph | steps: apply the relevant results; compute}]
An acidic indicator HIn ($K_a=10^{-9}$) has the ionised form of red colour and the unionised form of blue colour. The red colour is distinctly visible to a human eye when the ionised form is at least 2 times the unionized form whereas blue colour is distinctly visible to a human eye when the unionised form is at least, 10 times the ionised form. Various titrations are studied using this indicator. When colour change is visible during the titration, end point is observed for different titrations given in list-I. [Given: $\log2=0.3$]
Match the lists:
\begin{center}
\begin{tabular}{|c|p{3.8cm}|c|p{2.0cm}|}
\hline
\multicolumn{2}{|p{4.1cm}|}{\centering List-I (Titration)} & \multicolumn{2}{p{2.3cm}|}{\centering List-II (pH-values)} \\
\hline
(P) & 0.1 M 100 ml HCl (in the flask) vs 0.1 M NaOH (in the burette). & (1) & Equivalence point pH = 9. \\
\hline
(Q) & {\sloppy 0.2 M 100 ml CH$_3$COOH $K_a=10^{-5}$ (in the flask) vs 0.2 M NaOH (in the burette).\par} & (2) & End point pH = 7. \\
\hline
(R) & 0.2 M 100 ml NaOH (in the flask) vs 0.2 M CH$_3$COOH $K_a=10^{-5}$ (in the burette). & (3) & End point pH = 8. \\
\hline
(S) & 0.1 M 100 ml NaOH (in the flask) vs 0.1 M HCl (in the burette). & (4) & End point pH = 9.3. \\
\hline
 & & (5) & End point pH = 10. \\
\hline
\end{tabular}
\end{center}
\begin{qoptions}
\opt{A}{P$\to$4;Q$\to$1,4;R$\to$1,3;S$\to$3}
\opt{B}{P$\to$5;Q$\to$1,5;R$\to$1,3;S$\to$3}
\opt{C}{P$\to$2;Q$\to$2;R$\to$2;S$\to$2}
\opt{D}{P$\to$2;Q$\to$1;R$\to$1;S$\to$2}
\end{qoptions}
\end{bankq}
\begin{bankq}[id=SII12-P1-CHEM-S1-Q03, source={Score-II - Test 12 (04 Feb 2024)}, paper=1, qno=3, type=single, subtopic={Solubility product}, answer={C}, answer_src=key, difficulty=4, diff_parts={K2 S1 I1 C1 F0}, diff_note={concepts: standard problem in this chapter, solubility product | steps: apply the relevant results; compute}]
On decreasing the pH from 7 to 4, the solubility of a sparingly soluble salt (MX) of a weak acid (HX) increased from $10^{-4}$ mol $L^{-1}$ to $10^{-3}$ mol $L^{-1}$. The $pK_a$ of HX is:
\begin{qoptions}
\opt{A}{3}
\opt{B}{4}
\opt{C}{6}
\opt{D}{2}
\end{qoptions}
\end{bankq}
\begin{bankq}[id=SII13-P1-CHEM-S1-Q04, source={Score-II - Test 13 (06 Feb 2024)}, paper=1, qno=4, type=single, subtopic={Buffers}, answer={A}, answer_src=key, difficulty=4, diff_parts={K2 S1 I1 C1 F0}, diff_note={concepts: standard problem in this chapter, buffers | steps: apply the relevant results; compute}]
If 20 mL of 0.1 M NaOH is added to 30 mL of 0.2 M $CH_3COOH$ ($pK_a=4.74$), the pH of the resulting solution is :
\begin{qoptions}
\opt{A}{4.44}
\opt{B}{9.56}
\opt{C}{8.96}
\opt{D}{9.26}
\end{qoptions}
\end{bankq}
\begin{bankq}[id=SII13-P1-CHEM-S2-Q02, source={Score-II - Test 13 (06 Feb 2024)}, paper=1, qno=2, type=numerical, subtopic={pH}, answer={5}, answer_src=key, difficulty=8, diff_parts={K3 S1 I2 C2 F1}, diff_note={concepts: standard problem in this chapter, ph, apply the relevant results | steps: apply the relevant results; compute}]
3g of acetic acid is added to 250 mL of 0.1 M HCl and the solution made up to 500 mL. To 20 mL of this solution $\dfrac12$ mL of 5 M NaOH is added. The pH of the solution is \_\_\_\_. (Nearest integer)\\
\relax[Given: $pK_a$ of acetic acid = 4.75, molar mass of acetic acid = 60 g/mol, $\log3=0.4771$]\\
Neglect any changes in volume
\end{bankq}
\begin{bankpassage}[id=SII14-P1-CHEM-S3-G1-2, source={Score-II - Test 14 (06 Feb 2024)}, questions={1-2}]
Depending on the pH of the solution, $Mg^{2+}$ ions may hydrolyze to water soluble $MgOH^+$ species or precipitate as $Mg(OH)_2$.\\
$Mg(OH)_2(s)\rightleftharpoons Mg^{2+}(aq)+2OH^-(aq)$\hfill $K_{sp}=6\times10^{-12}$\\
$Mg^{2+}(aq)+OH^-(aq)\rightleftharpoons MgOH^+(aq)$\hfill $K_1=4\times10^2$\\
Consider a 10 litre solution of $2\times10^{-4}$ M $MgCl_2$ prepared in a buffer of pH = 11.\\
\relax[Atomic mass : Mg = 24, Cl = 35.5, O = 16, H = 1]
\end{bankpassage}
\begin{bankq}[id=SII14-P1-CHEM-S3-Q01, source={Score-II - Test 14 (06 Feb 2024)}, paper=1, qno=1, type=numerical, subtopic={Solubility product}, answer={71.45}, answer_src=key, difficulty=8, diff_parts={K3 S1 I2 C2 F1}, diff_note={concepts: standard problem in this chapter, solubility product, apply the relevant results | steps: apply the relevant results; compute}, group=SII14-P1-CHEM-S3-G1-2]
Determine \% $Mg^{2+}$ by mole (relative to all soluble Mg species) in this solution.
\end{bankq}
\begin{bankq}[id=SII14-P1-CHEM-S3-Q02, source={Score-II - Test 14 (06 Feb 2024)}, paper=1, qno=2, type=numerical, subtopic={Solubility product}, answer={111.12}, answer_src=solved, working={Official key reads "111.12 or 111…" (second value truncated); banked 111.12.}, difficulty=8, diff_parts={K3 S1 I2 C2 F1}, diff_note={concepts: standard problem in this chapter, solubility product, apply the relevant results | steps: apply the relevant results; compute}, group=SII14-P1-CHEM-S3-G1-2]
Calculate the mass of $Mg(OH)_2$ (in milli-grams) that will precipitate out from this solution.
\end{bankq}
\begin{bankq}[id=SII17-P1-CHEM-S1-Q04, source={Score-II - Test 17 (13 Feb 2024)}, paper=1, qno=4, type=single, subtopic={pH}, answer={C}, answer_src=solved, working={Ka=1e-4: initial pH 2.5 (A false); half-equivalence pH=4 but equivalence pH=8.41 not 9.09 (B false); weak acid-strong base: phenolphthalein suitable, methyl orange not (C).}, difficulty=4, diff_parts={K2 S1 I1 C1 F0}, diff_note={concepts: standard problem in this chapter, ph | steps: apply the relevant results; compute}]
20 ml, 0.1M HA $\left[(K_b)_{A^-}=10^{-10}\right]$ is titrated against 0.2M, NaOH. Select the correct statement :-\\
\relax[Given : log2 = 0.3, log3 = 0.48]
\begin{qoptions}
\opt{A}{Initially pH of HA solution is 11.5}
\opt{B}{At half equivalence point pH of solution is 4 and at equivalence point pH is 9.09}
\opt{C}{Phenolphthalein is suitable indicator for this titration but methyl orange is not suitable}
\opt{D}{Phenolphthalein \& methyl orange both are suitable indicator for end point determination.}
\end{qoptions}
\end{bankq}
\begin{bankq}[id=SII18-P1-CHEM-S2-Q07, source={Score-II - Test 18 (13 Feb 2024)}, paper=1, qno=7, type=multi, subtopic={pH}, answer={B,C,D}, answer_src=key, difficulty=8, diff_parts={K3 S1 I2 C2 F1}, diff_note={concepts: standard problem in this chapter, ph, apply the relevant results | steps: apply the relevant results; compute}]
The correct increasing order(s) of pH of 0.1 M -- aqueous solution of the following compounds, at $25^\circ$C, is/are

\begin{qoptions}
\opt{A}{HCl < CH$_3$COOH < HCOOH}
\opt{B}{H$_2$SO$_4$ < HCl < H$_3$PO$_4$}
\opt{C}{CH$_3$COOH < CH$_3$COONH$_4$ < CH$_3$COONa}
\opt{D}{NH$_4$Cl < CH$_3$COONH$_4$ < NH$_3$}
\end{qoptions}
\end{bankq}
\begin{bankq}[id=SII19-P1-CHEM-S2-Q06, source={Score-II - Test 19 (20 Feb 2024)}, paper=1, qno=6, type=numerical, subtopic={Solubility product}, answer={18}, answer_src=key, difficulty=8, diff_parts={K3 S1 I2 C2 F1}, diff_note={concepts: standard problem in this chapter, solubility product, apply the relevant results | steps: apply the relevant results; compute}]
The solubility product of Cr(OH)$_3$ at 298 K is $6.0\times10^{-31}$. The concentration of hydroxide ions in a saturated solution of Cr(OH)$_3$ will be $(x\times10^{-31})^{1/4}$, then x is :
\end{bankq}
\begin{bankq}[id=SII21-P1-CHEM-S3-Q03, source={Score-II - Test 21 (14 Mar 2024)}, paper=1, qno=3, type=numerical, subtopic={Buffers}, answer={Bonus}, answer_src=solved, working={Printed key omits or corrupts this numeric answer and the question was not solved; banked as Bonus.}, difficulty=8, diff_parts={K3 S1 I2 C2 F1}, diff_note={concepts: standard problem in this chapter, buffers, apply the relevant results | steps: apply the relevant results; compute}]
A solution is prepared by mixing 0.02 mol each of H$_2$CO$_3$, NaHCO$_3$, Na$_2$CO$_3$, and NaOH in 100 mL of water. pH of the resulting solution is\underline{\hspace{2em}}.
[Given : $pK_{a1}$ and $pK_{a2}$ of H$_2$CO$_3$ are 6.37 and 10.32, respectively ; $\log 2 = 0.30$]
\end{bankq}
\begin{bankq}[id=SII21-P2-CHEM-S3-Q01, source={Score-II - Test 21 (14 Mar 2024)}, paper=2, qno=1, type=numerical, subtopic={pH}, answer={6}, answer_src=key, difficulty=8, diff_parts={K3 S1 I2 C2 F1}, diff_note={concepts: standard problem in this chapter, ph, apply the relevant results | steps: apply the relevant results; compute}]
Concentration of H$_2$SO$_4$ and Na$_2$SO$_4$ in a solution is 1 M and $1.8\times10^{-2}$ M, respectively. Molar solubility of PbSO$_4$ in the same solution is $X\times10^{-Y}$ M (expressed in scientific notation). The value of Y is\underline{\hspace{2em}}.
[Given: Solubility product of PbSO$_4$ ($K_{sp}$) $=1.6\times10^{-8}$. For H$_2$SO$_4$, $K_{a1}$ is very large and $K_{a2}=1.2\times10^{-2}$]
\end{bankq}
\begin{bankq}[id=SII22-P1-CHEM-S1-Q04, source={Score-II - Test 22 (14 Mar 2024)}, paper=1, qno=4, type=single, subtopic={Solubility product}, answer={B}, answer_src=key, difficulty=4, diff_parts={K2 S1 I1 C1 F0}, diff_note={concepts: standard problem in this chapter, solubility product | steps: apply the relevant results; compute}]
The solubility of $PbF_2$ (formula weight = 245) is 0.46g/L. What is the solubility product :--
\begin{qoptions}
\opt{A}{$1.1\times10^{-10}$}
\opt{B}{$2.6\times10^{-8}$}
\opt{C}{$1.1\times10^{-7}$}
\opt{D}{$6.8\times10^{-9}$}
\end{qoptions}
\end{bankq}
\begin{bankq}[id=SII23-P1-CHEM-S3-Q02, source={Score-II - Test 23 (17 Mar 2024)}, paper=1, qno=2, type=numerical, subtopic={pH}, answer={1.75}, answer_src=key, difficulty=8, diff_parts={K3 S1 I2 C2 F1}, diff_note={concepts: standard problem in this chapter, ph, apply the relevant results | steps: apply the relevant results; compute}]
10 mL of H$_2$A (weak diprotic acid) solution is titrated against 0.1M NaOH. pH of the solution is plotted against volume of strong base added and following observation is made.
\bfig{90.4}{SII23_p1_chem_s3_q02_1.png}
If pH of the solution at first equivalence point is pH$_1$ and at second equivalence point is pH$_2$.
Calculate the value of (pH$_2-$pH$_1$) at $25^\circ$C
Given: For H$_2$A, $pK_{a_1}=4.6$ and $pK_{a_2}=8$, $\log25=1.4$
\end{bankq}
\begin{bankq}[id=SII24-P1-CHEM-S1-Q09, source={Score-II - Test 24 (17 Mar 2024)}, paper=1, qno=9, type=single, subtopic={Solubility product}, answer={C}, answer_src=key, difficulty=4, diff_parts={K2 S1 I1 C1 F0}, diff_note={concepts: standard problem in this chapter, solubility product | steps: apply the relevant results; compute}]
For a 1-2 type sparingly soluble electrolyte, solubility product is $4\times10^{-12}$. The concentration of cation in saturated solution will be :
\begin{qoptions}
\opt{A}{$2\times10^{-6}$ M}
\opt{B}{$10^{-4}$ M}
\opt{C}{$2\times10^{-4}$ M}
\opt{D}{$10^{-6}$ M}
\end{qoptions}
\end{bankq}
\begin{bankq}[id=SII24-P1-CHEM-S2-Q06, source={Score-II - Test 24 (17 Mar 2024)}, paper=1, qno=6, type=numerical, subtopic={pH}, answer={85}, answer_src=key, difficulty=8, diff_parts={K3 S1 I2 C2 F1}, diff_note={concepts: standard problem in this chapter, ph, apply the relevant results | steps: apply the relevant results; compute}]
If the pKa of lactic acid is 5, then the pH of 0.005 M calcium lactate solution at 25$^\circ$ C is \underline{\hspace{2em}} $\times10^{-1}$.
\end{bankq}
\begin{bankq}[id=SII26-P1-CHEM-S1-Q04, source={Score-II - Test 26 (20 Mar 2024)}, paper=1, qno=4, type=single, subtopic={Buffers}, answer={B}, answer_src=key, difficulty=4, diff_parts={K2 S1 I1 C1 F0}, diff_note={concepts: standard problem in this chapter, buffers | steps: apply the relevant results; compute}]
A buffer solution is prepared by mixing `X' moles of CH$_3$COONa and `y' moles of CH$_3$COOH such that (x + y) = 1, in water to make 1 L buffer solution. If the buffer capacity of this buffer solution is plotted against moles of salt CH$_3$COONa (with x) then the plot obtained will be (to scale approximately).
\begin{qoptions}
\opt{A}{\bfigi{105.3}{SII26_p1_chem_s1_q04_A.png}}
\opt{B}{\bfigi{103.8}{SII26_p1_chem_s1_q04_B.png}}
\opt{C}{\bfigi{103.6}{SII26_p1_chem_s1_q04_C.png}}
\opt{D}{\bfigi{104.1}{SII26_p1_chem_s1_q04_4.png}}
\end{qoptions}
\end{bankq}
\begin{bankq}[id=SII27-P1-CHEM-S2-Q02, source={Score-II - Test 27 (23 Mar 2024)}, paper=1, qno=2, type=numerical, subtopic={Solubility product}, answer={32}, answer_src=key, difficulty=8, diff_parts={K3 S1 I2 C2 F1}, diff_note={concepts: standard problem in this chapter, solubility product, apply the relevant results | steps: apply the relevant results; compute}]
When 1M Pb(NO$_3$)$_2$ solution is tirated with 1 M KI solution then what will be the osmotic pressure (in atm) when equivalence point was reached at 300K. (R = 0.08 atm-L/mole-K)
\end{bankq}
\begin{bankq}[id=SII28-P1-CHEM-S3-Q01, source={Score-II - Test 28 (23 Mar 2024)}, paper=1, qno=1, type=numerical, subtopic={Solubility product}, answer={13.20}, answer_src=key, difficulty=8, diff_parts={K3 S1 I2 C2 F1}, diff_note={concepts: standard problem in this chapter, solubility product, apply the relevant results | steps: apply the relevant results; compute}]
250 mL of 0.16 M MgCl$_2$(aq) is mixed with 750 mL of 0.2 M NH$_3$(aq) resulting in precipitation of some Mg(OH)$_2$. To redissolve Mg(OH)$_2$, what minimum mass (in gram) of (NH$_4$)$_2$SO$_4$ should be added to the above solution?\\
(Given : K$_{sp}$ [Mg(OH)$_2$] = $9\times10^{-12}$, K$_b$ (NH$_3$) = $2\times10^{-5}$)
\end{bankq}
\begin{bankq}[id=SII29-P1-CHEM-S1-Q02, source={Score-II - Test 29 (26 Mar 2024)}, paper=1, qno=2, type=single, subtopic={Solubility product}, answer={C}, answer_src=key, difficulty=4, diff_parts={K2 S1 I1 C1 F0}, diff_note={concepts: standard problem in this chapter, solubility product | steps: apply the relevant results; compute}]
A$_3$B$_2$ is a sparingly soluble salt of molar mass M (g mole$^{-1}$) and solubility x g lit$^{-1}$. The ratio of the molar concentration of B$^{3-}$ to the solubility product of the salt is

\begin{qoptions}
\opt{A}{$108\dfrac{x^5}{M^5}$}
\opt{B}{$\dfrac1{108}\dfrac{M^4}{x^4}$}
\opt{C}{$\dfrac1{54}\dfrac{M^4}{x^4}$}
\opt{D}{None}
\end{qoptions}
\end{bankq}
\begin{bankq}[id=SII30-P1-CHEM-S1-Q02, source={Score-II - Test 30 (27 Mar 2024)}, paper=1, qno=2, type=single, subtopic={Solubility product}, answer={B}, answer_src=key, difficulty=4, diff_parts={K2 S1 I1 C1 F0}, diff_note={concepts: standard problem in this chapter, solubility product | steps: apply the relevant results; compute}]
Solubility of calcium phosphate (molecular mass, $M$) in water is $W_g$ per 100 mL at 25$^\circ$ C. Its solubility product at 25$^\circ$C will be approximately.
\begin{qoptions}
\opt{A}{$10^7\left(\dfrac WM\right)^3$}
\opt{B}{$10^7\left(\dfrac WM\right)^5$}
\opt{C}{$10^3\left(\dfrac WM\right)^5$}
\opt{D}{$10^5\left(\dfrac WM\right)^5$}
\end{qoptions}
\end{bankq}
\begin{bankq}[id=SII31-P1-CHEM-S1-Q19, source={Score-II - Test 31 (29 Mar 2024)}, paper=1, qno=19, type=single, subtopic={pH}, answer={D}, answer_src=key, difficulty=4, diff_parts={K2 S1 I1 C1 F0}, diff_note={concepts: standard problem in this chapter, ph | steps: apply the relevant results; compute}]
Three solution of strong electrolytes 25 ml of 0.1 M HX, 25 ml of 0.1 M H$_2$Y and 50 ml of 0.1N Z(OH)$_2$ are mixed. What is the pOH of the resulting solution at 25$^\circ$C ? [Given log 5 = 0.7]
\begin{qoptions}
\opt{A}{1.6}
\opt{B}{7}
\opt{C}{11.6}
\opt{D}{12.4}
\end{qoptions}
\end{bankq}
\begin{bankq}[id=SII31-P1-CHEM-S2-Q10, source={Score-II - Test 31 (29 Mar 2024)}, paper=1, qno=10, type=single, subtopic={Buffers}, answer={9}, answer_src=key, difficulty=7, diff_parts={K3 S1 I2 C2 F0}, diff_note={concepts: standard problem in this chapter, buffers, apply the relevant results | steps: apply the relevant results; compute}]
What volume (in mL) of a 1.0 M HCl solution should be added to in aqueous solution of 1.0 L 0.001 M HCl in order to decrease its pH by one unit ?
\end{bankq}
\begin{bankq}[id=SII33-P1-CHEM-S1-Q01, source={Score-II - Test 33 (01 Apr 2024)}, paper=1, qno=1, type=single, subtopic={pH}, answer={Bonus}, answer_src=solved, working={Printed key block for Paper-1 is misaligned/truncated (duplicated and overlapping cells) and the question was not solved; banked as Bonus.}, difficulty=7, diff_parts={K3 S1 I2 C2 F0}, diff_note={concepts: standard problem in this chapter, ph, apply the relevant results | steps: apply the relevant results; compute}]
K$_{a_1}$ and K$_{a_2}$ are the 1$^{\text{st}}$ and 2$^{\text{nd}}$ ionization constants of maleic acid whereas K$'_{a_1}$ and K$'_{a_2}$ are those of fumaric acid, then choose the correct options among the following .
\begin{qoptions}
\opt{A}{$pK_{a_1}>pK'_{a_1}$}
\opt{B}{$pK_{a_2}>pK'_{a_2}$}
\opt{C}{$pK_{a_1}<pK'_{a_1}$}
\opt{D}{$pK_{a_2}<pK'_{a_2}$}
\end{qoptions}
\end{bankq}
\begin{bankq}[id=SII33-P1-CHEM-S2-Q01, source={Score-II - Test 33 (01 Apr 2024)}, paper=1, qno=1, type=numerical, subtopic={pH}, answer={Bonus}, answer_src=solved, working={Printed key block for Paper-1 is misaligned/truncated (duplicated and overlapping cells) and the question was not solved; banked as Bonus.}, difficulty=8, diff_parts={K3 S1 I2 C2 F1}, diff_note={concepts: standard problem in this chapter, ph, apply the relevant results | steps: apply the relevant results; compute}]
At 25$^\circ$C, the degree of dissociation of pure water is $1.9\times10^{-9}$. If the specific conductance of water at this temperature is x $\times10^{-6}$ Sm$^{-1}$. Calculate the value of x. Given ionic molar conductances at infinite dilution for H$^+$ and OH$^-$ ions are $350\times10^{-4}$ and $200\times10^{-4}$ Sm$^2$ mole$^{-1}$ respectively
\end{bankq}
\begin{bankq}[id=SII33-P2-CHEM-S1-Q01, source={Score-II - Test 33 (01 Apr 2024)}, paper=2, qno=1, type=multi, subtopic={pH}, answer={B,C,D}, answer_src=key, difficulty=8, diff_parts={K3 S1 I2 C2 F1}, diff_note={concepts: standard problem in this chapter, ph, apply the relevant results | steps: apply the relevant results; compute}]
$V_1$ ml of 3M HCl solution is mixed with $V_2$ ml of 3M H$_2$SO$_4$ to form a solution in which molarity of H$^+$ ions is 5 M, then the value of $V_1$ and $V_2$ may be
\begin{qoptions}
\opt{A}{150 ml, 75 ml}
\opt{B}{150 ml, 300 ml}
\opt{C}{125 ml, 250 ml}
\opt{D}{100 ml, 200 ml}
\end{qoptions}
\end{bankq}
\begin{bankq}[id=SII35-P1-CHEM-S1-Q01, source={Score-II - Test 35 (04 Feb 2025)}, paper=1, qno=1, type=single, subtopic={pH}, answer={C}, answer_src=key, difficulty=4, diff_parts={K2 S1 I1 C1 F0}, diff_note={concepts: standard problem in this chapter, ph | steps: apply the relevant results; compute}]
Concentration of H$_2$SO$_4$ and Na$_2$SO$_4$ in a solution is 1 M and $1.8\times10^{-2}$ M, respectively. Molar solubility of PbSO$_4$ in the same solution is :\\
\relax[Given: Solubility product of PbSO$_4$(K$_{sp}$) = $1.6\times10^{-8}$. For H$_2$SO$_4$, Ka$_1$ is very large and Ka$_2$ = $1.2\times10^{-4}$]
\begin{qoptions}
\opt{A}{$\dfrac89\times10^{-6}$ M}
\opt{B}{$\sqrt{1.6}\times10^{-4}$ M}
\opt{C}{$1.33\times10^{-4}$ M}
\opt{D}{$1.6\times10^{-6}$ M}
\end{qoptions}
\end{bankq}
\begin{bankq}[id=SII36-P1-CHEM-S2-Q09, source={Score-II - Test 36 (07 Feb 2025)}, paper=1, qno=9, type=multi, subtopic={pH}, answer={A,B,C}, answer_src=key, difficulty=8, diff_parts={K3 S1 I2 C2 F1}, diff_note={concepts: standard problem in this chapter, ph, apply the relevant results | steps: apply the relevant results; compute}]
In a 100 ml 0.1 M CH$_3$COOH solution, 0.1 M Ca(OH)$_2$ is slowly added. Then which one of the following statement(s) is/are correct (K$_a$ of CH$_3$COOH $=10^{-5}$, log 3 = 0.5, log 2 = 0.3)

\begin{qoptions}
\opt{A}{pH of the solution will be 5 after adding 25 ml of the base.}
\opt{B}{pH of the solution will be nearly 8.9 after adding 50 ml of the base.}
\opt{C}{Before adding the base, initial pH of the solution was nearly 3.}
\opt{D}{After adding 100 ml of the base pH of the solution becomes 13.}
\end{qoptions}
\end{bankq}
\begin{bankq}[id=SII38-P1-CHEM-S2-Q12, source={Score-II - Test 38 (12 Feb 2025)}, paper=1, qno=12, type=multi, subtopic={pH}, answer={A,B,C}, answer_src=key, difficulty=8, diff_parts={K3 S1 I2 C2 F1}, diff_note={concepts: standard problem in this chapter, ph, apply the relevant results | steps: apply the relevant results; compute}]
A 0.252 g sample of monoprotic organic base (amine) is dissolved in water and titrated against 0.07 M H$_2$SO$_4$ solution at 25$^\circ$C. After addition of 20 ml of acid, a pH of 10.7 is recorded. Equivalence point is observed when 40 ml of H$_2$SO$_4$ is added. (log 2 = 0.3, log 3 = 0.48, log 5 = 0.7)\\
Based on the above information select correct statement(s)
\begin{qoptions}
\opt{A}{Molar mass of the organic base is 45 g/mole}
\opt{B}{Resultant solution after addition of 20 ml of H$_2$SO$_4$ can act as basic buffer solution.}
\opt{C}{$K_b$ of the organic base is $5\times10^{-4}$}
\opt{D}{Solution formed at equivalence point is basic in nature.}
\end{qoptions}
\end{bankq}
\begin{bankq}[id=SII39-P1-CHEM-S3-Q02, source={Score-II - Test 39 (02 Mar 2025)}, paper=1, qno=2, type=single, subtopic={pH}, answer={5}, answer_src=key, difficulty=7, diff_parts={K3 S1 I2 C2 F0}, diff_note={concepts: standard problem in this chapter, ph, apply the relevant results | steps: apply the relevant results; compute}]
Calculate the pH of solution prepared by mixing 40 ml of 0.1M acetic acid with 20 ml of 0.1M sodium hydroxide. Given $K_{a(\text{ACETIC ACID})}=10^{-5}$
\end{bankq}
\begin{bankq}[id=SII40-P1-CHEM-S1-Q03, source={Score-II - Test 40 (16 Mar 2025)}, paper=1, qno=3, type=single, subtopic={Solubility product}, answer={A}, answer_src=key, difficulty=4, diff_parts={K2 S1 I1 C1 F0}, diff_note={concepts: standard problem in this chapter, solubility product | steps: apply the relevant results; compute}]
On decreasing the $p$H from 7 to 2, the solubility of a sparingly soluble salt (MX) of a weak acid (HX) increased from $10^{-4}$ mol L$^{-1}$ to `x' mol L$^{-1}$. Then `x' is ($K_a=10^{-4}$) :
\begin{qoptions}
\opt{A}{$10^{-3}$}
\opt{B}{$10^{-4}$}
\opt{C}{$10^{-5}$}
\opt{D}{$10^{-2}$}
\end{qoptions}
\end{bankq}
\begin{bankq}[id=SII41-P1-CHEM-S1-Q03, source={Score-II - Test 41 (19 Mar 2025)}, paper=1, qno=3, type=single, subtopic={Solubility product}, answer={C}, answer_src=key, difficulty=4, diff_parts={K2 S1 I1 C1 F0}, diff_note={concepts: standard problem in this chapter, solubility product | steps: apply the relevant results; compute}]
\textbf{Assertion :-} Solubility of Pb(OH)$_2$ (weak base) decreases on addition of NaOH.\\
\textbf{Reason :-} NaOH increases OH$^{-}$ concentration forcing the equilibrium to shift in the backward direction.
\begin{qoptions}
\opt{A}{Assertion and Reason both are correct and reason is correct explanation of assertion.}
\opt{B}{Assertion and Reason both are correct and reason is incorrect explanation of assertion.}
\opt{C}{Assertion and Reason both are incorrect.}
\opt{D}{Assertion is correct and Reason is incorrect.}
\end{qoptions}
\end{bankq}
\begin{bankq}[id=SII41-P1-CHEM-S1-Q05, source={Score-II - Test 41 (19 Mar 2025)}, paper=1, qno=5, type=single, subtopic={pH}, answer={A}, answer_src=key, difficulty=4, diff_parts={K2 S1 I1 C1 F0}, diff_note={concepts: standard problem in this chapter, ph | steps: apply the relevant results; compute}]
The aqueous solution of which of the following will have highest freezing point?
\begin{qoptions}
\opt{A}{Glucose (0.1 m)}
\opt{B}{NaCl (0.1 m)}
\opt{C}{Na$_3$PO$_4$ (0.05 m)}
\opt{D}{CH$_3$COOH (0.2 m, degree of dissociation = 0.1)}
\end{qoptions}
\end{bankq}
\begin{bankq}[id=SII42-P1-CHEM-S1-Q02, source={Score-II - Test 42 (19 Mar 2025)}, paper=1, qno=2, type=multi, subtopic={pH}, answer={A,B,C}, answer_src=key, difficulty=8, diff_parts={K3 S1 I2 C2 F1}, diff_note={concepts: standard problem in this chapter, ph, apply the relevant results | steps: apply the relevant results; compute}]
If pH of an aqueous solution of 0.01M monoprotic acid HA is found to be 3 at T kelvin, then the correct information(s) is/are

\begin{qoptions}
\opt{A}{Degree of dissociation of the acid would be 10\%}
\opt{B}{Osmotic pressure of the solution would be 0.011RT}
\opt{C}{Theoretical molar mass of HA is 1.1 times of its observed molar mass}
\opt{D}{K$_b$ of A$^-$(aq.) is $10^{-4}$ (approx.) at 298K}
\end{qoptions}
\end{bankq}
\begin{bankq}[id=SII42-P1-CHEM-S2-Q03, source={Score-II - Test 42 (19 Mar 2025)}, paper=1, qno=3, type=numerical, subtopic={Solubility product}, answer={99.98}, answer_src=key, difficulty=8, diff_parts={K3 S1 I2 C2 F1}, diff_note={concepts: standard problem in this chapter, solubility product, apply the relevant results | steps: apply the relevant results; compute}, also_in={Score Advanced - Test 17 (25 May 2023) P2 Q3; Score Advanced - Test 35 (15 May 2024) P2 Q3}]
Solid AgNO$_3$ is added to an aqueous solution containing 0.2 M Br$^-$ and 0.1 M I$^-$. Assuming that AgNO$_3$ addition does not change volume of this solution, calculate the percentage of [I$^-$] precipitated, when [Br$^-$] just starts precipitating: (Given: $k_{sp}$(AgBr) = $10^{-13}$ M$^2$; $k_{sp}$(AgI) = $10^{-17}$ M$^2$)
\end{bankq}
\begin{bankq}[id=SII42-P1-CHEM-S3-Q01, source={Score-II - Test 42 (19 Mar 2025)}, paper=1, qno=1, type=numerical, subtopic={pH}, answer={5}, answer_src=key, difficulty=8, diff_parts={K3 S1 I2 C2 F1}, diff_note={concepts: standard problem in this chapter, ph, apply the relevant results | steps: apply the relevant results; compute}]
A solution of weak acid HA was titrated with base NaOH. The equivalence point was reached when 50 mL of 0.1 M NaOH has been added. 25 mL of 0.1 M HCl was added to titrated solution and the pH was found to be 5.0. The pk$_a$ value of weak acid HA is \_\_\_\_\_.
\end{bankq}
\begin{bankq}[id=SII43-P1-CHEM-S1-Q04, source={Score-II - Test 43 (23 Mar 2025)}, paper=1, qno=4, type=single, subtopic={Hydrolysis}, answer={B}, answer_src=key, difficulty=4, diff_parts={K2 S1 I1 C1 F0}, diff_note={concepts: standard problem in this chapter, hydrolysis | steps: apply the relevant results; compute}]
When hydrated ferric chloride is dissolved in aqueous solution of oxalic acid containing potassium hydroxide
\begin{qoptions}
\opt{A}{Brown ppt of Fe(OH)$_3$ is obtained}
\opt{B}{Green crystals of K$_3$[Fe(C$_2$O$_4$)$_3$] are obtained}
\opt{C}{Yellow crystals of K[FeCl$_4$] are obtained}
\opt{D}{Yellow solution of [Fe(H$_2$O)$_6$]Cl$_3$ \& K$_2$C$_2$O$_4$ are obtained}
\end{qoptions}
\end{bankq}
\begin{bankq}[id=SII43-P1-CHEM-S1-Q09, source={Score-II - Test 43 (23 Mar 2025)}, paper=1, qno=9, type=single, subtopic={pH}, answer={B}, answer_src=key, difficulty=4, diff_parts={K2 S1 I1 C1 F0}, diff_note={concepts: standard problem in this chapter, ph | steps: apply the relevant results; compute}]
pH of the aq. solution of 0.05 M Ca(CN)$_2$ at 25$^\circ$C is :\\
$\left\{\mathrm{pK_b(CN^-)=8.26}\right\}$
\begin{qoptions}
\opt{A}{10.63}
\opt{B}{9.37}
\opt{C}{8.63}
\opt{D}{8.37}
\end{qoptions}
\end{bankq}
\begin{bankq}[id=SII44-P1-CHEM-S2-Q07, source={Score-II - Test 44 (23 Mar 2025)}, paper=1, qno=7, type=multi, subtopic={Solubility product}, answer={B,D}, answer_src=key, difficulty=8, diff_parts={K3 S1 I2 C2 F1}, diff_note={concepts: standard problem in this chapter, solubility product, apply the relevant results | steps: apply the relevant results; compute}]
0.1 mole AgNO$_3$ is added in 250m$\ell$ of saturated solution of AgC$\ell$ at 25$^\circ$C without changing the volume of solution.

Given: K$_{SP}$ of AgC$\ell$ $=1.0\times10^{-10}$ M$^2$ at 25$^\circ$C

Limiting molar conductance of Ag$^+$ ion $=60\ \Omega^{-1}$cm$^2$mol$^{-1}$

Limiting molar conductance of C$\ell^-$ ion $=75\ \Omega^{-1}$cm$^2$mol$^{-1}$

Limiting molar conductance of NO$_3^-$ ion $=75\ \Omega^{-1}$cm$^2$mol$^{-1}$

Which of the following statement(s) is/are \textbf{CORRECT}?

\begin{qoptions}
\opt{A}{O$_2$ gas is liberated at anode, if this solution is electrolysed by using inert electrodes}
\opt{B}{Total conductivity of this solution $\approx5.4\times10^{-2}\ \Omega^{-1}$cm$^{-1}$}
\opt{C}{Concentration of C$\ell^-$ in resulting solution is $2.5\times10^{-9}$ M}
\opt{D}{In the final solution, concentration of Ag$^+$ is 0.4 M}
\end{qoptions}
\end{bankq}
\begin{bankq}[id=SII45-P1-CHEM-S1-Q15, source={Score-II - Test 45 (26 Mar 2025)}, paper=1, qno=15, type=single, subtopic={pH}, answer={C}, answer_src=key, difficulty=4, diff_parts={K2 S1 I1 C1 F0}, diff_note={concepts: standard problem in this chapter, ph | steps: apply the relevant results; compute}]
A solution at 25$^\text{o}$C is 0.01 M in acetic acid and 0.05 M in sodium acetate. Then the pOH of solution is.\\
$\big(K_a(\text{acetic acid})=1.8\times10^{-5};$\\ $\log5=0.7;\ \log2=0.3;\ \log3=0.47\ \&\ \log3.47=0.54\big)$
\begin{qoptions}
\opt{A}{4.45}
\opt{B}{5.45}
\opt{C}{8.55}
\opt{D}{9.55}
\end{qoptions}
\end{bankq}
\begin{bankq}[id=SII47-P1-CHEM-S1-Q10, source={Score-II - Test 47 (30 Mar 2025)}, paper=1, qno=10, type=single, subtopic={pH}, answer={C}, answer_src=key, difficulty=4, diff_parts={K2 S1 I1 C1 F0}, diff_note={concepts: standard problem in this chapter, ph | steps: apply the relevant results; compute}]
The moles of H$^+$ from H$_2$O alone in a $1l,\ \sqrt{5}\times10^{-7}\text{ M HCl}$ solution at $25^\circ$C is $(\sqrt5=2.23)$
\begin{qoptions}
\opt{A}{$10^{-7}$}
\opt{B}{$6.85\times10^{-8}$}
\opt{C}{$3.85\times10^{-8}$}
\opt{D}{$10^{-8}$}
\end{qoptions}
\end{bankq}

\section*{Score Advanced papers}
\begin{bankq}[id=SA02-P1-CHEM-S1-Q01, source={Score Advanced - Test 02 (13 May 2021)}, paper=1, qno=1, type=multi, subtopic={pH}, answer={B,C}, answer_src=solved, working={Dilution lowers [H+]; neutral water at higher temperature; 1e-4 M HCl has pH 4 against pH 3 for 0.1 M HCN}, difficulty=4, diff_parts={K1 S2 I1 C0 F1}, diff_note={concepts: acidity statements | steps: dilution; temperature; compare pH}]
Which statement(s) is/are correct -
\begin{qoptions}
\opt{A}{If 0.1 M CH$_3$COOH aqueous solution (K$_a$ = 1.8 $\times$ 10$^{-5}$) is diluted at 25$^\circ$C, then [H$^+$] will increase.}
\opt{B}{On heating a sample of pure water [H$^+$] and [OH$^-$] increases but the sample still remains neutral}
\opt{C}{10$^{-4}$ M HCl solution is less acidic than 0.1 M HCN (K$_{a\,\mathrm{HCN}}$ = 10$^{-5}$) solution}
\opt{D}{10$^{-8}$ M HCl aqueous solution will have pH exactly equal to 8}
\end{qoptions}
\end{bankq}
\begin{bankq}[id=SA02-P2-CHEM-S1-Q01, source={Score Advanced - Test 02 (13 May 2021)}, paper=2, qno=1, type=multi, subtopic={Hydrolysis}, answer={A,B,C}, answer_src=solved, working={Kb=1e-5 gives [OH-]=1e-3, pH 11, [HCO3-]=1e-3, [H2CO3]=1e-9; CO3 2- is about 0.1}, difficulty=7, diff_parts={K2 S2 I2 C1 F1}, diff_note={concepts: carbonate hydrolysis, hydrolysis equilibrium | steps: hydrolysis equilibrium; concentrations of species; pH}]
For 0.1 M aq. solution of Na$_2$CO$_3$ at 25$^\circ$C, select the correct options (Use : H$_2$CO$_3$ : K$_{a_1}$ = 10$^{-5}$, K$_{a_2}$ = 10$^{-9}$ ) and K$_w$ (H$_2$O) = 10$^{-14}$
\begin{qoptions}
\opt{A}{pH of the solution is 11.}
\opt{B}{Equilibrium concentration of H$_2$CO$_3$ is 10$^{-9}$ M.}
\opt{C}{Equilibrium concentration of HCO$_3^-$ is 10$^{-3}$ M}
\opt{D}{Equilibrium concentration of CO$_3^{2-}$ is 0.01 M}
\end{qoptions}
\end{bankq}
\begin{bankq}[id=SA04-P1-CHEM-S2-Q01, source={Score Advanced - Test 04 (04 Jul 2021)}, paper=1, qno=1, type=multi, subtopic={pH}, answer={C,D}, answer_src=solved, working={Curve 2 is HCl and 1 is a weak acid with Ka=(10\^{}-3)\^{}2/0.1=10\^{}-5. The end point of the weak acid is basic so phenolphthalein works}, difficulty=7, diff_parts={K2 S2 I2 C1 F1}, diff_note={concepts: titration curves of two acids, ph | steps: initial pH of 0.1 M: curve 2 starts at 1 (strong), curve 1 at 3 (weak); Ka=10\^{}-5; indicator at the end point}]
With respect to above information in paragraph indentify correct statement(s)
\begin{qoptions}
\opt{A}{Curve (1) may be of HCl while (2) may be of CH$_3$COOH.}
\opt{B}{The nature of solution at equivalence point for curve(1) is acidic}
\opt{C}{Phenolphthalein can be used to detect the equivalence point for curve (1)}
\opt{D}{For weaker acid among (1) and (2) the value of K$_\text{a}$ = 10$^{-5}$}
\end{qoptions}
\end{bankq}
\begin{bankq}[id=SA04-P1-CHEM-S2-Q02, source={Score Advanced - Test 04 (04 Jul 2021)}, paper=1, qno=2, type=multi, subtopic={Buffers}, answer={A,C}, answer_src=solved, working={10\^{}0.6=3.98 which is above 0.25; pink is visible; unionised fraction is 20 percent}, difficulty=6, diff_parts={K2 S1 I2 C1 F1}, diff_note={concepts: indicator colour at a given pH, buffers | steps: ratio ionised/unionised = 10\^{}(pH-pKa)}]
Let the colour of an acidic indicator HIn (colourless) be visible only when it's ionized form (pink) is 25\% or more than the unionized form in a solution. If HIn (pK$_\text{a}$ = 9) is added to solution of pH = 9.6, then correct observation(s) is/are -
\begin{qoptions}
\opt{A}{Pink colour will be visible}
\opt{B}{Pink colour will not be visible}
\opt{C}{Unionized form of indicator is less than 25\% in solution}
\opt{D}{Ionized form of indicator is less than 25\% in solution}
\end{qoptions}
\end{bankq}
\begin{bankq}[id=SA04-P2-CHEM-S1-Q01, source={Score Advanced - Test 04 (04 Jul 2021)}, paper=2, qno=1, type=single, subtopic={Buffers}, answer={D}, answer_src=solved, working={P leaves excess NaOH (no); Q leaves 1 mmol acetate with 0.5 mmol acid; R is NH3 with NH4Cl; S is F- with HF; so Q, R, S}, difficulty=4, diff_parts={K2 S1 I1 C1 F0}, diff_note={concepts: which mixtures are buffers, moles of base and acid after neutralisation | steps: moles of base and acid after neutralisation}]
Observe the following four solutions :\par
P : 10 ml, 0.1 M NaOH + 5ml, 0.1 M HCl\par
Q : 10 ml, 0.1 M NaOH + 15ml, 0.1 M CH$_3$COOH\par
R : 10 ml, 0.1 M NH$_3$ + 10ml, 0.1 M NH$_4$Cl\par
S : 10 ml, 0.05 M NaF + 5ml, 0.1 M HF\par
Which of above solutions act as buffer
\begin{qoptions}
\opt{A}{P and Q only}
\opt{B}{R and S only}
\opt{C}{P, R, S only}
\opt{D}{Q, R, S only}
\end{qoptions}
\end{bankq}
\begin{bankq}[id=SA04-P2-CHEM-S2-Q03, source={Score Advanced - Test 04 (04 Jul 2021)}, paper=2, qno=3, type=numerical, subtopic={Solubility product}, answer={1.2}, answer_src=solved, working={1.992e-2/332=6e-5 M, [Ag+]=1.2e-4 M, X=1.2}, difficulty=5, diff_parts={K2 S1 I1 C1 F1}, diff_note={concepts: silver ion concentration from solubility, solubility product | steps: solubility in mol per litre; two Ag+ per formula unit}]
The solubility of silver chromate is 1.992 $\times$ 10$^{-2}$ gm /litre. The molar mass of Ag$_2$CrO$_4$ is 332 gm/mol. If the molar concentration of Ag$^+$ in saturated solution of Ag$_2$CrO$_4$ is X $\times$ 10$^{-4}$ M, give `X'.
\end{bankq}
\begin{bankq}[id=SA05-P1-CHEM-S2-Q08, source={Score Advanced - Test 05 (15 Jul 2021)}, paper=1, qno=8, type=numerical, subtopic={Solubility product}, answer={-3}, answer_src=key, difficulty=5, diff_parts={K2 S1 I1 C1 F1}, diff_note={concepts: solubility of Al(OH)3 in a buffer, solubility product | steps: [OH-]=10\^{}-10; S=Ksp/[OH-]\^{}3}]
If the solubility of Al(OH)$_3$, (K$_{\text{sp}}$ = 10$^{-33}$) in a buffer solution of pH = 4 is 10$^x$. Give x?
\end{bankq}
\begin{bankq}[id=SA06-P2-CHEM-S3-Q02, source={Score Advanced - Test 06 (04 Sep 2021)}, paper=2, qno=2, type=numerical, subtopic={Buffers}, answer={8}, answer_src=key, difficulty=6, diff_parts={K2 S1 I2 C1 F1}, diff_note={concepts: titration of sodium sesquicarbonate, buffers | steps: phenolphthalein end point with carbonate to bicarbonate; methyl orange end point}]
40 ml of 0.05 M solution of sodium sesquicarbonate (Na$_2$CO$_3$.NaHCO$_3$. 2H$_2$O) is titrated against 0.05 M HCl. ``X'' ml of HCl is used when phenolphthalein is the indicator and ``Y'' ml of HCl is used when methyl orange is the indicator in two separate titrations. Hence, value of $\dfrac{(Y-X)}{10}$ is
\end{bankq}
\begin{bankq}[id=SA07-P1-CHEM-S3-Q04, source={Score Advanced - Test 07 (16 Apr 2023)}, paper=1, qno=4, type=numerical, subtopic={Solubility product}, answer={10}, answer_src=key, difficulty=8, diff_parts={K3 S1 I2 C2 F1}, diff_note={concepts: solubility ratio in a buffer to pure water, solubility product, solubility in water from Ksp | steps: solubility in water from Ksp; in buffer from Ka}]
The ionisation constant of benzoic acid is 6.46 $\times$ 10$^{-6}$ and K$_\text{sp}$ for silver benzoate is 2.5 $\times$ 10$^{-13}$. The number of times silver benzoate is more soluble in a buffer of pH 3.19 compared to its solubility in pure water will be .........\par
[Given : Antilog [--3.19] = 6.46 $\times$ 10$^{-4}$ M]
\end{bankq}
\begin{bankq}[id=SA09-P1-CHEM-S2-Q09, source={Score Advanced - Test 09 (29 Apr 2023)}, paper=1, qno=9, type=multi, subtopic={Solubility product}, answer={A,B}, answer_src=key, difficulty=9, diff_parts={K3 S2 I2 C2 F1}, diff_note={concepts: silver bromide in aqueous ammonia, solubility product, complex formation | steps: complex formation; concentrations from Ksp and Kf; conductivities}]
NH$_3$ concentration is a saturated solution of AgBr(s) is maintained at 0.2M. The correct information(s) related with the solution is/are :\par
K$_\text{sp}$ of Ag Br = 2 $\times$ 10$^{-10}$\par
K$_\text{f}$ of Ag(NH$_3$)$_2^+$ = 8 $\times$ 10$^6$\par
$\lambda^\text{o}_{\text{(Ag}^+)}$ = 6 $\times$ 10$^{-3}$ Sm$^2$ mol$^{-1}$\par
$\lambda^\text{o}_{\text{(Br}^-)}$ = 8 $\times$ 10$^{-3}$ Sm$^2$ mol$^{-1}$
\begin{qoptions}
\opt{A}{Conductivity of Br$^-$ = 6.4 $\times$ 10$^{-2}$ S m$^{-1}$}
\opt{B}{Conductivity of Ag$^+$ = 1.5 $\times$ 10$^{-7}$ S m$^{-1}$}
\opt{C}{[Ag$^+$] = 8 $\times$ 10$^{-3}$M}
\opt{D}{[Br$^-$] = 2 $\times$ 10$^{-3}$M}
\end{qoptions}
\end{bankq}
\begin{bankq}[id=SA10-P1-CHEM-S3-Q05, source={Score Advanced - Test 10 (06 May 2023)}, paper=1, qno=5, type=numerical, subtopic={Buffers}, answer={4}, answer_src=key, difficulty=6, diff_parts={K2 S1 I2 C1 F1}, diff_note={concepts: buffer pH change when salt is added, buffers | steps: Henderson equation; change of 0.6 in pH means a factor of 4}]
To a solution of acetic acid, solid sodium acetate is added gradually forming a buffer. When x mole of salt is added, the pH has a certain value. If y mole of salt is added, the pH at this time changes by 0.6 units. Find the ratio (y/x). Given that y $>$ x.\par
[log2 = 0.3, log3 = 0.48, log5 = 0.7.]
\end{bankq}
\begin{bankq}[id=SA10-P2-CHEM-S2-Q08, source={Score Advanced - Test 10 (06 May 2023)}, paper=2, qno=8, type=multi, subtopic={Solubility product}, answer={A,B,D}, answer_src=key, difficulty=8, diff_parts={K3 S1 I2 C2 F1}, diff_note={concepts: precipitation of AgI before AgCl, solubility product, compare the silver ion concentration needed | steps: compare the silver ion concentration needed}]
To an aqueous solution containing a mixture of 0.05M NaCl and 0.05M NaI when AgNO$_3$ is slowly added to such a solution, then correct option(s) is/are :\par
(K$_\text{sp}$AgCl = 10$^{-10}$, K$_\text{sp}$AgI = 4 $\times$ 10$^{-16}$)
\begin{qoptions}
\opt{A}{Concentration of Ag$^+$ required to precipitate Cl$^\ominus$ is 2 $\times$ 10$^{-9}$M}
\opt{B}{Concentration of Ag$^+$ required to precipitate I$^\ominus$ is 8 $\times$ 10$^{-15}$M}
\opt{C}{AgCl and AgI will precipitate out of the solution together}
\opt{D}{First AgI will precipitate}
\end{qoptions}
\end{bankq}
\begin{bankq}[id=SA12-P1-CHEM-S1-Q03, source={Score Advanced - Test 12 (11 May 2023)}, paper=1, qno=3, type=single, subtopic={Solubility product}, answer={B}, answer_src=key, difficulty=7, diff_parts={K3 S1 I2 C2 F0}, diff_note={concepts: concentration of free Cu+ in an excess cyanide solution, solubility product, complex formation equilibrium | steps: complex formation equilibrium}, also_in={Score Advanced - Test 29 (01 May 2024) P2 Q3}]
A 0.1 M solution of Cu$^+$ ions is stirred with an excess of potassium cyanide sufficient to convert all the Cu$^+$ ions to the corresponding cuprocyanide complex [Cu(CN)$_4$]$^{-3}$ and in addition to providing the solution with an excess of CN$^-$ equal to 0.2 M. Calculate the maximum pH of the solution when the final solution is treated with hydrogen sulphide to maintain [H$_2$S] = 0.1M and such that the precipitation of cuprous sulphide is prevented. The instability constant for [Cu(CN)$_4$]$^{-3}$ is 6.4 $\times$ 10$^{-15}$ M$^4$, K$_\text{a, overall}$ of H$_2$S = 1.6 $\times$ 10$^{-21}$ M, K$_\text{sp}$ of Cu$_2$S = 2.56 $\times$ 10$^{-27}$M$^3$.
\begin{qoptions}
\opt{A}{4.0}
\opt{B}{10.0}
\opt{C}{10.8}
\opt{D}{3.2}
\end{qoptions}
\end{bankq}
\begin{bankq}[id=SA13-P1-CHEM-S1-Q07, source={Score Advanced - Test 13 (13 May 2023)}, paper=1, qno=7, type=numerical, subtopic={pH}, answer={4.8}, answer_src=key, difficulty=6, diff_parts={K2 S1 I2 C1 F1}, diff_note={concepts: pH of a hydrolysing iron(II) nitrate solution, h=sqrt(Kh/C)... [H+]=sqrt(Kh C) | steps: h=sqrt(Kh/C)... [H+]=sqrt(Kh C)}]
The pH of 0.05 M-Fe(NO$_3$)$_2$ solution is (K$_\text{h}$ of Fe$^{2+}$ = 5.0 $\times$ 10$^{-9}$, log 2 = 0.30).\par
\end{bankq}
\begin{bankpassage}[id=SA13-P2-CHEM-S2-G7-8, source={Score Advanced - Test 13 (13 May 2023)}, questions={7-8}]
In water, HF behaves as a weak acid. The reactions involved are :\par
\textbf{Reaction-I :} HF(aq) + H$_2$O($\ell$) $\rightleftharpoons$ H$_3$O$^+$(aq) + F$^-$(aq); K$_1$ = 1.1 $\times$ 10$^{-3}$\par
\textbf{Reaction-II :} HF(aq) + F$^-$(aq) $\rightleftharpoons$ HF$_2^-$(aq); K$_2$ = 0.26\par
Two chemists wanted to determine the acidity constant, K$_1$, using the same HF solution and measuring the pH of solution. One chemist was having the idea of reaction-I only while the other chemist was having the idea of both reactions. Surprisingly, both the chemist determined the same value of K$_1$ (as given above).\par
\end{bankpassage}
\begin{bankq}[id=SA13-P2-CHEM-S2-Q07, source={Score Advanced - Test 13 (13 May 2023)}, paper=2, qno=7, type=single, subtopic={pH}, answer={A}, answer_src=key, difficulty=7, diff_parts={K3 S1 I2 C2 F0}, diff_note={concepts: hydrofluoric acid equilibrium with its dimer form, ph, combine K1 and K2 | steps: combine K1 and K2; pH of the solution}, group=SA13-P2-CHEM-S2-G7-8]
Which of the following must be incorrect ?
\begin{qoptions}
\opt{A}{[HF]$_0$ = $\sqrt{2\text{K}_1\text{K}_2}$}
\opt{B}{K$_1$ = [H$_3$O$^+$]}
\opt{C}{[HF]$_0$ = 2K$_1$}
\opt{D}{[HF]$_0$ = 2[H$_3$O$^+$]}
\end{qoptions}
\end{bankq}
\begin{bankq}[id=SA13-P2-CHEM-S2-Q08, source={Score Advanced - Test 13 (13 May 2023)}, paper=2, qno=8, type=single, subtopic={pH}, answer={B}, answer_src=key, difficulty=7, diff_parts={K3 S1 I2 C2 F0}, diff_note={concepts: initial concentration of HF, ph, from the equilibrium constants and the measured pH | steps: from the equilibrium constants and the measured pH}, group=SA13-P2-CHEM-S2-G7-8]
The initial concentration of HF solution used is :
\begin{qoptions}
\opt{A}{0.22 M}
\opt{B}{2.2 $\times$ 10$^{-3}$ M}
\opt{C}{3.3 $\times$ 10$^{-3}$ M}
\opt{D}{2.0 M}
\end{qoptions}
\end{bankq}
\begin{bankq}[id=SA15-P1-CHEM-S3-Q02, source={Score Advanced - Test 15 (18 May 2023)}, paper=1, qno=2, type=numerical, subtopic={Buffers}, answer={4.5}, answer_src=key, difficulty=8, diff_parts={K3 S1 I2 C2 F1}, diff_note={concepts: pH of a mixture of HCl, acetate and formate, buffers, HCl converts acetate to acetic acid | steps: HCl converts acetate to acetic acid; equilibrium with formate with K=0.1}, also_in={Score Advanced - Test 32 (08 May 2024) P1 Q2}]
Calculate the pH of solution formed by addition of 1 mole each of HCl, CH$_3$COONa and HCOONa in enough water to make 1 L solution. The pK$_a$ of HCOOH and CH$_3$COOH are 4 and 5, respectively.
\end{bankq}
\begin{bankpassage}[id=SA16-P1-CHEM-S3-G1-2, source={Score Advanced - Test 16 (20 May 2023)}, questions={1-2}]
100 ml 0.1 M weak acid HA (K$_a$ = 10$^{-6}$) is titrated with V ml 0.1 M NaOH (aq).\par
[Given : log 3 = 0.48, log 2 = 0.30]\par
\end{bankpassage}
\begin{bankq}[id=SA16-P1-CHEM-S3-Q01, source={Score Advanced - Test 16 (20 May 2023)}, paper=1, qno=1, type=numerical, subtopic={Buffers}, answer={5.52}, answer_src=key, difficulty=7, diff_parts={K3 S1 I2 C1 F1}, diff_note={concepts: pH of a partially neutralised weak acid, buffers, Henderson equation with 1:3 ratio | steps: Henderson equation with 1:3 ratio}, group=SA16-P1-CHEM-S3-G1-2]
The pH of final solution when V = 25 ml is
\end{bankq}
\begin{bankq}[id=SA16-P1-CHEM-S3-Q02, source={Score Advanced - Test 16 (20 May 2023)}, paper=1, qno=2, type=numerical, subtopic={Buffers}, answer={9.35}, answer_src=key, difficulty=7, diff_parts={K3 S1 I2 C1 F1}, diff_note={concepts: pH at the equivalence point, hydrolysis, salt hydrolysis | steps: salt hydrolysis; pH=7+pKa/2+log c/2}, group=SA16-P1-CHEM-S3-G1-2]
The pH of final solution when V = 100 ml is
\end{bankq}
\begin{bankq}[id=SA17-P1-CHEM-S1-Q04, source={Score Advanced - Test 17 (25 May 2023)}, paper=1, qno=4, type=multi, subtopic={pH}, answer={A,C,D}, answer_src=key, difficulty=9, diff_parts={K3 S2 I2 C2 F1}, diff_note={concepts: dissociation of a weak monoprotic acid from its pH, Ka from pH, degree of dissociation | steps: Ka from pH; degree of dissociation; effect of temperature}, also_in={Score Advanced - Test 35 (15 May 2024) P1 Q4}]
If pH of an aqueous solution of 0.1M monoprotic acid HA is found to be 3 at T kelvin then
\begin{qoptions}
\opt{A}{Degree of dissociation of the acid would be 1 \% (approx.)}
\opt{B}{Osmotic pressure of the solution would be 0.011RT}
\opt{C}{Theoretical molar mass of HA is approximately 1.01 times of its observed molar mass}
\opt{D}{K$_b$ of A$^-_{(aq.)}$ is 10$^{-11}$ M (approx.) at 298K}
\end{qoptions}
\end{bankq}
\begin{bankq}[id=SA19-P1-CHEM-S3-Q03, source={Score Advanced - Test 19 (04 Apr 2024)}, paper=1, qno=3, type=numerical, subtopic={Solubility product}, answer={3}, answer_src=key, difficulty=8, diff_parts={K3 S1 I2 C2 F1}, diff_note={concepts: saturated salt solution with boiling point elevation, solubility product, molality from colligative data | steps: molality from colligative data}]
The saturated solution of a certain salt show elevation in boiling point equal to $3.12\times10^{-3}$ K in aqueous solution. If the solubility of salt is $2\times10^{-3}$ mole/kg at boiling point, then how many moles of ions are given by 1 mole of solid salt on complete dissociation. ($K_b$(H$_2$O) $=0.52$ K kg/mole)
\end{bankq}
\begin{bankpassage}[id=SA23-P1-CHEM-S3-G1-2, source={Score Advanced - Test 23 (14 Apr 2024)}, questions={1-2}]
Salt hydrolysis : This is the reverse reaction of acid base neutralisation\par
Salt of many types :\par
(i) Salt of strong acid and strong base\par
(ii) Salt of strong base and weak acid\par
(iii) Salt of strong acid and weak base\par
(iv) Salt of weak acid and weak base\par
Answer the following questions
\end{bankpassage}
\begin{bankq}[id=SA23-P1-CHEM-S3-Q01, source={Score Advanced - Test 23 (14 Apr 2024)}, paper=1, qno=1, type=numerical, subtopic={Hydrolysis}, answer={7.40}, answer_src=key, difficulty=8, diff_parts={K3 S1 I2 C2 F1}, diff_note={concepts: pH of a mixture of sodium carbonate and hydrochloric acid, hydrolysis, reaction gives bicarbonate | steps: reaction gives bicarbonate; amphiprotic pH}, group=SA23-P1-CHEM-S3-G1-2]
0.2 M aq. solution of Na$_2$CO$_3$ is mixed with 0.2 M aq. HCl in equal volume proportion at 25$^\circ$C. pH of the resulting solution will be\par
(For H$_2$CO$_3$: K$_{a_1}$ = 2 $\times$ 10$^{-6}$, K$_{a_2}$ = 8 $\times$ 10$^{-10}$, log (2) = 0.3)
\end{bankq}
\begin{bankq}[id=SA23-P1-CHEM-S3-Q02, source={Score Advanced - Test 23 (14 Apr 2024)}, paper=1, qno=2, type=numerical, subtopic={Hydrolysis}, answer={0.02}, answer_src=key, difficulty=8, diff_parts={K3 S1 I2 C2 F1}, diff_note={concepts: carbonate concentration in the final solution, hydrolysis, disproportionation of the amphiprotic species | steps: disproportionation of the amphiprotic species}, group=SA23-P1-CHEM-S3-G1-2]
In the above final solution, if concentration of CO$_3^{2-}$ is X molar, then value of 10 X will be
\end{bankq}
\begin{bankq}[id=SA25-P2-CHEM-S4-Q05, source={Score Advanced - Test 25 (21 Apr 2024)}, paper=2, qno=5, type=single, subtopic={Solubility product}, answer={300}, answer_src=key, difficulty=7, diff_parts={K3 S1 I2 C2 F0}, diff_note={concepts: solubility product from a logarithmic plot, read intercepts of the plot | steps: read intercepts of the plot}, also_in={Score Advanced - Test 26 (21 Apr 2024) P2 Q5}]
A metal M form sparingly soluble metal sulphate. The logarithm of concentration of dissolved metal ion in a solution saturated in the metal sulphate is plotted below as a function of logarithm of concentration of sulphate ion in solution. The molecular formula of metal sulphate is M$_x$(SO$_4$)$_y$. Then find z $\times$ 10$^{-6}$.\par
$z=\dfrac{x+y}{K_{sp}\text{ of M}_x\text{(SO}_4)_y}$, x and y are minimum integer co-efficient.\par
\bfig{133.3}{SA25_p2_chem_s4_q05_1.png}
\end{bankq}
\begin{bankq}[id=SA27-P2-CHEM-S1-Q08, source={Score Advanced - Test 27 (24 Apr 2024)}, paper=2, qno=8, type=multi, subtopic={Buffers}, answer={B,C,D}, answer_src=solved, working={1 mmol of each. Phenolphthalein needs 1 mmol HCl = 10 mL of 0.1 M HCl (A false, 1 M is wrong). From pp to methyl orange: 2 mmol = 20 mL of 0.1 M (B true). Methyl orange from the start: 3 mmol = 6 mL of 0.5 M (C true). Equivalent weight of Na2CO3 is 106 and 53 (D true).}, difficulty=8, diff_parts={K3 S1 I2 C2 F1}, diff_note={concepts: titration of a bicarbonate carbonate mixture with two indicators, buffers, phenolphthalein end point converts CO3 2- to HCO3- | steps: phenolphthalein end point converts CO3 2- to HCO3-; methyl orange converts all to H2CO3}]
10 ml mixture containing 0.1 M NaHCO$_3$ and 0.1M Na$_2$CO$_3$ is titrated against an acid with an acid$-$base indicator. Among the following which observations is/are correct
\begin{qoptions}
\opt{A}{When phenolphthalene indicator used, end point observed when 10ml of 1M HCl is added}
\opt{B}{After getting end point with phenolphthalein if we added methyl orange indicator, to get next end point 20 ml of 0.1M HCl is to be added furtherly}
\opt{C}{When methyl orange indicator used from the beginning, end point observed when 6ml of 0.5M HCl added}
\opt{D}{When phenolphthalene indicator used, gram equivalent weight of Na$_2$CO$_3$ is 106 g but when methyl orange indicator used gram equivalent weight of Na$_2$ CO$_3$ is 53g}
\end{qoptions}
\end{bankq}
\begin{bankq}[id=SA27-P2-CHEM-S3-Q05, source={Score Advanced - Test 27 (24 Apr 2024)}, paper=2, qno=5, type=numerical, subtopic={Solubility product}, answer={7}, answer_src=solved, working={Essentially all silver is [Ag(NH3)2]+ = 1e-3 M. [Ag(NH3)+] = K1 [Ag(NH3)2+]/[NH3] = 1.4e-6 x 1e-3/0.2 = 7.0e-9 M. The mantissa P = 7 (the form 'P x 10\^{}-P' in the statement is inconsistent; 7.0e-9 is the computed value).}, difficulty=8, diff_parts={K3 S1 I2 C2 F1}, diff_note={concepts: concentration of the monoammine silver complex, solubility product, stepwise dissociation constants | steps: stepwise dissociation constants; free NH3 about 0.2 M}]
In the mixture 1.0 $\times$ 10$^{-3}$ M AgNO$_3$ + 0.20 M NH$_3$, concentration of [Ag (NH$_3$)]$^+$ in correct scientific notation is P $\times$ 10$^{-\mathrm P}$M. The value of `P' is\par
Given : [Ag(NH$_3$)$_2$]$^+$ $\rightleftharpoons$ [Ag(NH$_3$)]$^+$ + NH$_3$; K$_1$ = 1.4 $\times$ 10$^{-6}$M\par
[Ag(NH$_3$)]$^+$ $\rightleftharpoons$ Ag$^+$ + NH$_3$ ; K$_2$ = 4.3 $\times$ 10$^{-4}$M
\end{bankq}
\begin{bankq}[id=SA28-P2-CHEM-S3-Q06, source={Score Advanced - Test 28 (27 Apr 2024)}, paper=2, qno=6, type=matching, subtopic={pH}, answer={P$\to$2; Q$\to$3; R$\to$4; S$\to$1}, answer_src=solved, working={(P) 1 mmol of base and 1 mmol of HCl give a 0.1 M salt with Kh = 1e-4 so [H+] = sqrt(Kh C) and pH = 2.5 (2). (Q) Na3A with 2.5 mol HCl leaves 0.5 mol H3A and 0.5 mol H2A-, pH = pKa1 = 3 (3). (R) alpha = 10/11 gives [A-]/[HA] = 10 and pH = 5+1 = 6 (4). (S) 1e-3 M HA with Ka = 1e-7 gives [H+] about 1e-5 so pH 5 (1).}, difficulty=9, diff_parts={K3 S2 I2 C2 F1}, diff_note={concepts: pH of solutions in various titration and dilution situations, hydrolysis of salt, buffer with equal acid and conjugate base | steps: hydrolysis of salt; buffer with equal acid and conjugate base; Henderson equation; weak acid dissociation}]
Match the Lists.\par
{\footnotesize\setlength{\tabcolsep}{1.5pt}\noindent\begin{tabular}{|c|L{4.9cm}|c|L{1.5cm}|}
\hline
 & \multicolumn{1}{>{\centering\arraybackslash}p{4.9cm}|}{List-I} & & \multicolumn{1}{>{\centering\arraybackslash}p{1.5cm}|}{List-II (pH)} \\ \hline
(P) & 2.5 ml $\dfrac25$ M weak base (K$_\mathrm b$ = 10$^{-10}$) is mixed with $\dfrac2{15}$ M HCl for complete neutralisation & (1) & 5 \\ \hline
(Q) & 1 mole Na$_3$A is mixed with 2.5 mole of HCl to give 2 L solution (pKa$_1$, pKa$_2$, pKa$_3$ of H$_3$A are 3, 6, 10 respectively) & (2) & 2.5 \\ \hline
(R) & pH at which 0.1 M CH$_3$COOH dissociates $\dfrac{1000}{11}$\% (pKa of CH$_3$COOH = 5) & (3) & 3 \\ \hline
(S) & 1 L 0.1 M HA is mixed with 99 L water (pKa of HA = pKb of A$^-$) & (4) & 6 \\ \hline
\end{tabular}}\par
\end{bankq}
\begin{bankq}[id=SA30-P2-CHEM-S3-Q06, source={Score Advanced - Test 30 (04 May 2024)}, paper=2, qno=6, type=numerical, subtopic={Solubility product}, answer={20}, answer_src=key, difficulty=8, diff_parts={K3 S1 I2 C2 F1}, diff_note={concepts: standard problem in this chapter, solubility product, apply the relevant results | steps: apply the relevant results; compute}, also_in={Score Advanced - Test 31 (04 May 2024) P2 Q6}]
The solubility of a salt (AB) of weak acid at pH = 3 is $\sqrt{\mathrm Y}\times10^{-3}$ mol L$^{-1}$. The value of Y is\underline{\hspace{0.8cm}}.\par
(Given that the value of solubility product of AB (K$_\mathrm{sp}$) = 2 $\times$ 10$^{-10}$ and the value of ionization constant of HB(K$_\mathrm a$) = 1 $\times$ 10$^{-8}$)
\end{bankq}
\begin{bankq}[id=SA33-P1-CHEM-S1-Q03, source={Score Advanced - Test 33 (12 May 2024)}, paper=1, qno=3, type=single, subtopic={Solubility product}, answer={D}, answer_src=key, difficulty=4, diff_parts={K2 S1 I1 C1 F0}, diff_note={concepts: standard problem in this chapter, solubility product | steps: apply the relevant results; compute}, also_in={Score Advanced - Test 37 (19 May 2024) P1 Q3}]
A small amount of solid lead iodide was added to a beaker of water, which was stirred. Most of the solid settled at the bottom of the beaker, but a little dissolved, establishing the equilibrium\par
PbI$_2$ (s) $\rightleftharpoons$ Pb$^{2+}$ (aq) + 2I$^-$ (aq)\par
The rates of the forward and reverse reactions were monitored over time, producing the graph shown below:\par
\bfig{165.7}{SA33_p1_chem_s1_q03_1.png}
What happened at time t?
\begin{qoptions}
\opt{A}{The beaker was cooled in an ice-bath.}
\opt{B}{A small amount of solid KI was added to the beaker.}
\opt{C}{A small amount of solid PbI$_2$ was removed from the beaker.}
\opt{D}{A small amount of water was added to the beaker.}
\end{qoptions}
\end{bankq}
\begin{bankq}[id=SA33-P2-CHEM-S2-Q07, source={Score Advanced - Test 33 (12 May 2024)}, paper=2, qno=7, type=multi, subtopic={Buffers}, answer={D}, answer_src=key, difficulty=5, diff_parts={K2 S1 I1 C1 F1}, diff_note={concepts: standard problem in this chapter, buffers | steps: apply the relevant results; compute}, also_in={Score Advanced - Test 34 (12 May 2024) P2 Q7}]
Select the correct statement(s) from the following.
\begin{qoptions}
\opt{A}{If the critical micelle concentration (CMC) of a soap is 10$^{-3}$ M, then 10$^{-4}$ M solution of this soap will have colloidal nature}
\opt{B}{Charge on the colloidal particles may be determined with the help of dialysis.}
\opt{C}{Peptization is the process of conversion of a colloidal sol into precipitate by adding some electrolyte.}
\opt{D}{When some oil-soluble dye is added in the water in oil emulsion, the background becomes coloured.}
\end{qoptions}
\end{bankq}
\begin{bankq}[id=SA33-P2-CHEM-S3-Q02, source={Score Advanced - Test 33 (12 May 2024)}, paper=2, qno=2, type=numerical, subtopic={Solubility product}, answer={2}, answer_src=key, difficulty=8, diff_parts={K3 S1 I2 C2 F1}, diff_note={concepts: standard problem in this chapter, solubility product, apply the relevant results | steps: apply the relevant results; compute}, also_in={Score Advanced - Test 34 (12 May 2024) P2 Q2}]
Cd$^{2+}$ ion is precipitate after passing H$_2$S in acidified solution. In the final solution 0.1 milli-mole of CdSO$_4$ are present in 10 ml acidic saturated solution of 0.10M H$_2$S. What is the pH of the solution after filtering off precipitate, boiling off H$_2$S and making the solution 100 ml by adding H$_2$O :\par
[K$_\mathrm{sp}$ of CdS = 10$^{-28}$, For H$_2$S : K$_{\mathrm a_1}$ = 10$^{-9}$, K$_{\mathrm a_2}$ = 10$^{-20}$]
\end{bankq}
\begin{bankq}[id=SA39-P1-CHEM-S1-Q02, source={Score Advanced - Test 39 (20 Apr 2025)}, paper=1, qno=2, type=single, subtopic={Buffers}, answer={C}, answer_src=key, difficulty=4, diff_parts={K2 S1 I1 C1 F0}, diff_note={concepts: standard problem in this chapter, buffers | steps: apply the relevant results; compute}]
A solution of weak acid HA was titrated with strong base KOH. The weak acid is completely neutralized by adding 60 mL of 0.1 M KOH. Now, 20mL of 0.1 M HCl was added to titrated solution, the pH was found to be 5.3. What will be the pH of solution obtained by mixing 50 mL of 0.1 M KOH and 50 mL of 0.1 M HA? (log 2= 0.3)
\begin{qoptions}
\opt{A}{8.00}
\opt{B}{8.7}
\opt{C}{8.85}
\opt{D}{9.00}
\end{qoptions}
\end{bankq}
\begin{bankq}[id=SA39-P2-CHEM-S1-Q06, source={Score Advanced - Test 39 (20 Apr 2025)}, paper=2, qno=6, type=multi, subtopic={Buffers}, answer={A,B,C}, answer_src=solved, working={Solved by hand (key missing): result A,B,C.}, difficulty=8, diff_parts={K3 S1 I2 C2 F1}, diff_note={concepts: standard problem in this chapter, buffers, apply the relevant results | steps: apply the relevant results; compute}]
pH of which of the following solution in close to 5 ? (pka of CH$_3$COOH = 5, K$_\mathrm a$ of NH$_4^+$ = 5.6 $\times$ 10$^{-10}$)
\begin{qoptions}
\opt{A}{50 mL of 0.01 M CH$_3$COOH + 50 mL of 0.01 M CH$_3$COONa}
\opt{B}{0.1 M solution of NH$_4$Cl}
\opt{C}{0.1 M weak monoprotic acid having K$_\mathrm a$ = 1 $\times$ 10$^{-9}$}
\opt{D}{5 $\times$ 10$^{-6}$ M Ba(OH)$_2$}
\end{qoptions}
\end{bankq}
\begin{bankq}[id=SA40-P1-CHEM-S1-Q02, source={Score Advanced - Test 40 (27 Apr 2025)}, paper=1, qno=2, type=multi, subtopic={Buffers}, answer={A,B,D}, answer_src=key, difficulty=8, diff_parts={K3 S1 I2 C2 F1}, diff_note={concepts: standard problem in this chapter, buffers, apply the relevant results | steps: apply the relevant results; compute}]
How many of these solutions (of equal volume) on mixing create a Buffer solution?
\begin{qoptions}
\opt{A}{0.1 M NaOH, 0.3 M HCN}
\opt{B}{0.3 M H$_2$SO$_4$, 0.8 M CH$_3$COO$^-$Na$^+$}
\opt{C}{0.2 M Ba(OH)$_2$, 0.4 M HCN}
\opt{D}{0.3 M H$_2$CO$_3$, 0.1 M NaOH}
\end{qoptions}
\end{bankq}
\begin{bankq}[id=SA41-P1-CHEM-S1-Q02, source={Score Advanced - Test 41}, paper=1, qno=2, type=single, subtopic={pH}, answer={D}, answer_src=key, difficulty=4, diff_parts={K2 S1 I1 C1 F0}, diff_note={concepts: standard problem in this chapter, ph | steps: apply the relevant results; compute}]
2.5 mL of $\dfrac25$ M weak monoacidic base (K$_\mathrm b$ = 1 $\times$ 10$^{-12}$ at 25$^\circ$C) is titrated with $\dfrac2{15}$ M HCl in water at 25$^\circ$C. The concentration of H$^+$ at equivalence point is\par
(K$_\mathrm w$ = 1 $\times$ 10$^{-14}$ at 25$^\circ$C)
\begin{qoptions}
\opt{A}{3.7 $\times$ 10$^{-13}$ M}
\opt{B}{3.2 $\times$ 10$^{-7}$ M}
\opt{C}{3.2 $\times$ 10$^{-2}$ M}
\opt{D}{2.7 $\times$ 10$^{-2}$ M}
\end{qoptions}
\end{bankq}
\begin{bankq}[id=SA41-P2-CHEM-S4-Q05, source={Score Advanced - Test 41}, paper=2, qno=5, type=numerical, subtopic={pH}, answer={1380}, answer_src=key, difficulty=8, diff_parts={K3 S1 I2 C2 F1}, diff_note={concepts: standard problem in this chapter, ph, apply the relevant results | steps: apply the relevant results; compute}]
At 298 K, pure T$_2$O has pT (like pH) is 7.60. Find out the pT of a solution prepared by adding 10 ml of 0.2 M TCl to 15 ml of 0.20 M NaOT. (Given log 2 = 0.30)\par
Report your answer by multiplying actual answer with 100.
\end{bankq}

\section*{SRGBT papers}
\begin{bankq}[id=SRGBT05-P1-CHEM-S1-Q21, source={SRGBT - Test 05 (23 Aug 2021)}, paper=1, qno=21, type=multi, subtopic={Solubility product}, answer={D}, answer_src=key, difficulty=5, diff_parts={K2 S1 I1 C1 F1}, diff_note={concepts: standard problem in this chapter, solubility product | steps: apply the relevant results; compute}]
An aqueous solution is 0.3M in Al$^{3+}$ and 3M in Mg$^{2+}$.\par
Select the only \textbf{INCORRECT} option\par
Given : (K$_\mathrm{sp}$)$_{\mathrm{Al(OH)}_3}$ = 3 $\times$ 10$^{-34}$, (K$_\mathrm{sp}$)$_{\mathrm{Mg(OH)}_2}$ = 3 $\times$ 10$^{-12}$ M$^3$
\begin{qoptions}
\opt{A}{At pH = 9 both Al(OH)$_3$ and Mg(OH)$_2$ will precipitate}
\opt{B}{At pOH = 12 neither Al(OH)$_3$ nor Mg(OH)$_2$ will precipitate}
\opt{C}{For selective precipitation of Al(OH)$_3$ only, pH of solution should be 3 to 8.}
\opt{D}{Mg(OH)$_2$ will precipitate at maximum pH = 8}
\end{qoptions}
\end{bankq}
\begin{bankq}[id=SRGBT05-P1-CHEM-S1-Q23, source={SRGBT - Test 05 (23 Aug 2021)}, paper=1, qno=23, type=multi, subtopic={pH}, answer={A}, answer_src=key, difficulty=5, diff_parts={K2 S1 I1 C1 F1}, diff_note={concepts: standard problem in this chapter, ph | steps: apply the relevant results; compute}]
40 ml, 0.35 M NaOH solution is added to 50 ml 0.6 N H$_3$PO$_4$ solution. The pH of the mixture will be:\par
(Ka$_1$, Ka$_2$ and Ka$_3$ are 10$^{-3}$, 10$^{-8}$, 10$^{-12}$ respectively).
\begin{qoptions}
\opt{A}{7.82}
\opt{B}{5.76}
\opt{C}{8.84}
\opt{D}{9.32}
\end{qoptions}
\end{bankq}
\begin{bankq}[id=SRGBT05-P1-CHEM-S1-Q25, source={SRGBT - Test 05 (23 Aug 2021)}, paper=1, qno=25, type=multi, subtopic={Hydrolysis}, answer={A,B,C,D}, answer_src=key, difficulty=8, diff_parts={K3 S1 I2 C2 F1}, diff_note={concepts: standard problem in this chapter, hydrolysis, apply the relevant results | steps: apply the relevant results; compute}]
Which of the following statement are correct when equal volumes of 0.02 M HOCl \& 0.2 M CH$_3$COOH are mixed ?\par
K$_{\mathrm a_1}$ (HOCl) = 2$\times$10$^{-4}$; K$_{\mathrm a_2}$ (CH$_3$COOH) = 2$\times$10$^{-5}$ ; log 2 = 0.3
\begin{qoptions}
\opt{A}{$\alpha_1$ = 10$^{-1}$}
\opt{B}{$\alpha_2$ = 10$^{-2}$}
\opt{C}{pH = 2.7}
\opt{D}{[H$^+$] = 2 $\times$ 10$^{-3}$}
\end{qoptions}
\end{bankq}
\begin{bankq}[id=SRGBT05-P1-CHEM-S1-Q27, source={SRGBT - Test 05 (23 Aug 2021)}, paper=1, qno=27, type=multi, subtopic={Buffers}, answer={A,B}, answer_src=key, difficulty=8, diff_parts={K3 S1 I2 C2 F1}, diff_note={concepts: standard problem in this chapter, buffers, apply the relevant results | steps: apply the relevant results; compute}]
Buffer capacity is defined in terms of buffer index which is the change in the concentration of buffer acid (or base) required for changing its pH value by one, keeping [C$_\mathrm{acid}$ + C$_\mathrm{salt}$] or [C$_\mathrm{base}$ + C$_\mathrm{salt}$] constant. The correct statements is/are for acidic buffer of HA and A$^-$
\begin{qoptions}
\opt{A}{Buffer capacity is highest when [HA] \& [A$^-$] are large and [HA] = [A$^-$]}
\opt{B}{Buffer capacity is highest when the pH is equal to (or very near) the pKa of the HA used.}
\opt{C}{For making a buffer solution with a pH = 2, among acetic acid [pKa = 4.74] \& chlorine acid [pKa = 1.95], acetic acid is better choice.}
\opt{D}{Among buffer A (0.1 M CH$_3$COOH \& 0.1 M CH$_3$COONa) and buffer B [0.2 M CH$_3$COOH \& 0.2 M CH$_3$COONa] buffer A is better buffer}
\end{qoptions}
\end{bankq}
\begin{bankq}[id=SRGBT05-P1-CHEM-S1-Q29, source={SRGBT - Test 05 (23 Aug 2021)}, paper=1, qno=29, type=multi, subtopic={pH}, answer={B,C}, answer_src=key, difficulty=8, diff_parts={K3 S1 I2 C2 F1}, diff_note={concepts: standard problem in this chapter, ph, apply the relevant results | steps: apply the relevant results; compute}]
A 1 gram pill contains morphine (M = 285.34), a monobasic compound whose conjugate acid has pKa = 8.21, in addition to an unknown amount of inert material. To analyze the morphine content of the pill, it is dissolved in 50 ml 0.1 M aq. HCl and the unreacted acid is back titrated with 0.1 M NaOH solution.\par
(I) If the pill contains only morphine, the titration will require 35 ml NaOH solution to reach the equivalence point. [(285.34)$^{-1}$ = 0.0035]\par
(II) Methyl red which changes colour from red to yellow between pH = 4.4 and pH = 6.2, would be a suitable indication for this titration, mentioned in statement(I).
\begin{qoptions}
\opt{A}{I is correct statement}
\opt{B}{I is incorrect statement}
\opt{C}{II is correct statement}
\opt{D}{II is incorrect}
\end{qoptions}
\end{bankq}
\begin{bankq}[id=SRGBT05-P1-CHEM-S1-Q30, source={SRGBT - Test 05 (23 Aug 2021)}, paper=1, qno=30, type=multi, subtopic={pH}, answer={A,B}, answer_src=key, difficulty=8, diff_parts={K3 S1 I2 C2 F1}, diff_note={concepts: standard problem in this chapter, ph, apply the relevant results | steps: apply the relevant results; compute}]
Which of the following statements are not correct ?
\begin{qoptions}
\opt{A}{The pH of pure water increases with increase of temperature.}
\opt{B}{The solubility of KAl(SO$_4$)$_2$ in terms of its solubility product is (K$_\mathrm{sp}$/2)$^{1/4}$.}
\opt{C}{The pH of a solution obtained by mixing 50 mL of 0.4 M HCl and 50 mL of 0.2 M NaOH would be one.}
\opt{D}{The standard dissociation constant of NH$_4$OH is 1.8 $\times$ 10$^{-5}$. The standard hydrolysis constant of NH$_4^+$ ions at 25$^\circ$C would be 5.56 $\times$ 10$^{-10}$.}
\end{qoptions}
\end{bankq}
\begin{bankpassage}[id=SRGBT05-P1-CHEM-S1-G33-34, source={SRGBT - Test 05 (23 Aug 2021)}, questions={33-34}]
Titration of diprotic acid (H$_2$A) by strong base has been summarised in following graph\par
\bfig{86.8}{SRGBT05_p1_chem_s0_q32_1.png}
\end{bankpassage}
\begin{bankq}[id=SRGBT05-P1-CHEM-S1-Q33, source={SRGBT - Test 05 (23 Aug 2021)}, paper=1, qno=33, type=multi, subtopic={Buffers}, answer={C}, answer_src=key, difficulty=5, diff_parts={K2 S1 I1 C1 F1}, diff_note={concepts: standard problem in this chapter, buffers | steps: apply the relevant results; compute}, group=SRGBT05-P1-CHEM-S1-G33-34]
What are the pH of solution when it shows maximum buffer action ?
\begin{qoptions}
\opt{A}{2.34, 6}
\opt{B}{6, 9.66}
\opt{C}{2.34, 9.66}
\opt{D}{6, 12}
\end{qoptions}
\end{bankq}
\begin{bankq}[id=SRGBT05-P1-CHEM-S1-Q34, source={SRGBT - Test 05 (23 Aug 2021)}, paper=1, qno=34, type=multi, subtopic={Buffers}, answer={B}, answer_src=key, difficulty=5, diff_parts={K2 S1 I1 C1 F1}, diff_note={concepts: standard problem in this chapter, ph | steps: apply the relevant results; compute}, group=SRGBT05-P1-CHEM-S1-G33-34]
What is the pH at II equivalence point?
\begin{qoptions}
\opt{A}{6.0}
\opt{B}{11.83}
\opt{C}{9.66}
\opt{D}{12}
\end{qoptions}
\end{bankq}
\begin{bankq}[id=SRGBT06-P1-CHEM-S1-Q36, source={SRGBT - Test 06 (29 Aug 2021)}, paper=1, qno=36, type=multi, subtopic={Solubility product}, answer={A,B,D}, answer_src=key, difficulty=8, diff_parts={K3 S1 I2 C2 F1}, diff_note={concepts: standard problem in this chapter, solubility product, apply the relevant results | steps: apply the relevant results; compute}]
Which of the following is/are correct.
\begin{qoptions}
\opt{A}{Ag$_2$S is much less soluble than Ag$_2$O.}
\opt{B}{Fe(OH)$_3$ is much less soluble than Fe(OH)$_2$.}
\opt{C}{BeCO$_3$ is less soluble than MgCO$_3$.}
\opt{D}{Melting point of AlCl$_3$ is much less than that of NaCl.}
\end{qoptions}
\end{bankq}
\begin{bankq}[id=SRGBT12-P1-CHEM-S1-Q19, source={SRGBT - Test 12 (10 Oct 2021)}, paper=1, qno=19, type=multi, subtopic={Solubility product}, answer={C}, answer_src=key, difficulty=5, diff_parts={K2 S1 I1 C1 F1}, diff_note={concepts: standard problem in this chapter, solubility product | steps: apply the relevant results; compute}]
The solubility of ethylamine in a solution with pH = 8.5 is (K$_\mathrm{sp}$ of ethylamine = 1.78 $\times$ 10$^{-10}$)\par
\textbf{(Given:} $\sqrt{10}=3.16$ \textbf{)}
\begin{qoptions}
\opt{A}{5.6 $\times$ 10$^{-3}$ M}
\opt{B}{3.8 $\times$ 10$^{-4}$ M}
\opt{C}{5.6 $\times$ 10$^{-5}$ M}
\opt{D}{3.16 $\times$ 10$^{-6}$ M}
\end{qoptions}
\end{bankq}
\begin{bankq}[id=SRGBT12-P1-CHEM-S1-Q21, source={SRGBT - Test 12 (10 Oct 2021)}, paper=1, qno=21, type=multi, subtopic={Solubility product}, answer={C}, answer_src=key, difficulty=5, diff_parts={K2 S1 I1 C1 F1}, diff_note={concepts: standard problem in this chapter, solubility product | steps: apply the relevant results; compute}]
In the excess of NH$_3$ (aq.), Cu$^{2+}$ ions forms a deep blue complex ion [Cu(NH$_3$)$_4$]$^{2+}$ which has K$_\mathrm f$ = 10$^{11}$ the concentration of Cu$^{2+}$ ion in solution prepared by adding 5 $\times$ 10$^{-3}$ mole of CuSO$_4$ to 0.5 L of 0.4 M NH$_3$ (aq) is
\begin{qoptions}
\opt{A}{$\dfrac{5\times10^{-5}}{(36)^4}$}
\opt{B}{$\dfrac{10^{-6}}{(36)^4}$}
\opt{C}{$\dfrac{10^{-5}}{(36)^4}$}
\opt{D}{$\dfrac{10^{-2}}{(39)^4}$}
\end{qoptions}
\end{bankq}
\begin{bankq}[id=SRGBT12-P1-CHEM-S1-Q23, source={SRGBT - Test 12 (10 Oct 2021)}, paper=1, qno=23, type=multi, subtopic={pH}, answer={C}, answer_src=key, difficulty=5, diff_parts={K2 S1 I1 C1 F1}, diff_note={concepts: standard problem in this chapter, ph | steps: apply the relevant results; compute}]
The indicator constant for an acidic indicator, HIn, is 5$\times$10$^{-6}$ M. This indicator appears only in the colour of acidic form when $\dfrac{[\mathrm{In}^-]}{[\mathrm{HIn}]}\le\dfrac1{20}$ and it appears only in the colour of basic form when $\dfrac{[\mathrm{HIn}]}{[\mathrm{In}^-]}\le\dfrac1{40}$. The pH range of indicator is
\begin{qoptions}
\opt{A}{4.3 -- 6.3}
\opt{B}{4.0 -- 6.6}
\opt{C}{4.0 -- 6.9}
\opt{D}{3.7 -- 6.6}
\end{qoptions}
\end{bankq}
\begin{bankq}[id=SRGBT12-P1-CHEM-S1-Q25, source={SRGBT - Test 12 (10 Oct 2021)}, paper=1, qno=25, type=multi, subtopic={pH}, answer={B,C}, answer_src=key, difficulty=8, diff_parts={K3 S1 I2 C2 F1}, diff_note={concepts: standard problem in this chapter, ph, apply the relevant results | steps: apply the relevant results; compute}]
\bfig{228.5}{SRGBT12_p1_chem_s0_q25_1.png}
If valves are opened sequentially (A to D) \& each valve lets out 20 mL of stock solution, then select correct statements (Temp = 25$^\circ$C)
\begin{qoptions}
\opt{A}{After opening (A), buffer is created in beaker "P".}
\opt{B}{After opening (B), buffer is created beaker "Q".}
\opt{C}{After opening (C), contents in beaker Q can still act as buffer even though solutions pH decreases.}
\opt{D}{After opening (D), solution of pH $<$ 7 is created in beaker "R".}
\end{qoptions}
\end{bankq}
\begin{bankq}[id=SRGBT12-P1-CHEM-S1-Q27, source={SRGBT - Test 12 (10 Oct 2021)}, paper=1, qno=27, type=multi, subtopic={Solubility product}, answer={B,C,D}, answer_src=key, difficulty=8, diff_parts={K3 S1 I2 C2 F1}, diff_note={concepts: standard problem in this chapter, solubility product, apply the relevant results | steps: apply the relevant results; compute}]
A solution contains 0.05 M each of NaCl and Na$_2$CrO$_4$. Solid AgNO$_3$ is gradually added to it. Which of the following facts is/are incorrect?\par
Given: K$_\mathrm{sp}$(AgCl) = 1.7 x 10$^{-10}$ M$^2$ and K$_\mathrm{sp}$(Ag$_2$CrO$_4$) = 1.9 x 10$^{-12}$ M$^3$.
\begin{qoptions}
\opt{A}{Cl$^-$ ions are precipitated first}
\opt{B}{CrO$_4^{2-}$ ions are precipitated first}
\opt{C}{Both Cl$^-$ and CrO$_4^{2-}$ ions are precipitated simultaneously}
\opt{D}{The second ion starts precipitating when half of the first ion is precipitated}
\end{qoptions}
\end{bankq}
\begin{bankq}[id=SRGBT12-P1-CHEM-S1-Q28, source={SRGBT - Test 12 (10 Oct 2021)}, paper=1, qno=28, type=multi, subtopic={Hydrolysis}, answer={A,B}, answer_src=key, difficulty=8, diff_parts={K3 S1 I2 C2 F1}, diff_note={concepts: standard problem in this chapter, hydrolysis, apply the relevant results | steps: apply the relevant results; compute}]
CN$^-_{\mathrm{(aq.)}}$ + H$_2$O $\rightleftharpoons$ HCN$_{\mathrm{(aq.)}}$ + OH$^-_{\mathrm{(aq.)}}$\par
Addition of which of the following reagent can make the reaction proceed towards forward direction.
\begin{qoptions}
\opt{A}{H$_2$O}
\opt{B}{HCl}
\opt{C}{NaBr}
\opt{D}{KOH}
\end{qoptions}
\end{bankq}
\begin{bankq}[id=SRGBT12-P1-CHEM-S1-Q29, source={SRGBT - Test 12 (10 Oct 2021)}, paper=1, qno=29, type=multi, subtopic={pH}, answer={A,B,C}, answer_src=key, difficulty=8, diff_parts={K3 S1 I2 C2 F1}, diff_note={concepts: standard problem in this chapter, ph, apply the relevant results | steps: apply the relevant results; compute}]
1 M solution of a dibasic acid H$_2$A $\left(\mathrm K_{\mathrm a_1}=10^{-6}\text{ and }\mathrm K_{\mathrm a_2}=10^{-10}\right)$ is titrated against 1 M NaOH solution. Choose the \textbf{correct} statements regarding the relative concentration of H$_2$A, HA$^-$ and A$^{2-}$ at various pH during titration.
\begin{qoptions}
\opt{A}{At pH = 3 ; [H$_2$A] is maximum.}
\opt{B}{At pH = 6.5 ; [HA$^-$] is maximum.}
\opt{C}{At pH = 10 ; [H$_2$A] is minimum.}
\opt{D}{At pH = 8 ; [A$^{2-}$] is maximum.}
\end{qoptions}
\end{bankq}
\begin{bankq}[id=SRGBT12-P1-CHEM-S1-Q31, source={SRGBT - Test 12 (10 Oct 2021)}, paper=1, qno=31, type=multi, subtopic={Solubility product}, answer={A,B,C}, answer_src=key, difficulty=8, diff_parts={K3 S1 I2 C2 F1}, diff_note={concepts: standard problem in this chapter, solubility product, apply the relevant results | steps: apply the relevant results; compute}]
Solubility product of AgCl (s) in water at 25$^\circ$C is \ 12 $\times$ 10$^{-10}$ M$^2$ then select the correct options.
\begin{qoptions}
\opt{A}{Solubility of AgCl in presence of \ 10$^{-5}$ M AgNO$_3$ is \ 3 $\times$ 10$^{-5}$ moles/L.}
\opt{B}{Solubility of AgCl in presence of excess of KCl will increase.}
\opt{C}{Solubility of AgCl in presence of 2M NH$_3$ is greater than that in pure water.}
\opt{D}{Solubility of AgCl remains unaffected on adding any amount of NaBr because there is no common ion effect.}
\end{qoptions}
\end{bankq}
\begin{bankpassage}[id=SRGBT12-P1-CHEM-S1-G33-34, source={SRGBT - Test 12 (10 Oct 2021)}, questions={33-34}]
Consider the following experimental set-up with stop cocks: S$_1$, S$_2$, S$_3$, S$_4$ and taps : T$_1$, T$_2$, T$_3$, T$_4$\par
\bfig{170.7}{SRGBT12_p1_chem_s0_q32_1.png}
\end{bankpassage}
\begin{bankq}[id=SRGBT12-P1-CHEM-S1-Q33, source={SRGBT - Test 12 (10 Oct 2021)}, paper=1, qno=33, type=multi, subtopic={pH}, answer={C}, answer_src=key, difficulty=5, diff_parts={K2 S1 I1 C1 F1}, diff_note={concepts: standard problem in this chapter, ph | steps: apply the relevant results; compute}, group=SRGBT12-P1-CHEM-S1-G33-34]
In which of the following cases will the solution in tub have a pH different from others:
\begin{qoptions}
\opt{A}{T$_1$ and T$_4$ are opened}
\opt{B}{only T$_2$ is opened}
\opt{C}{T$_1$, T$_2$, T$_3$ and T$_4$ are all opened}
\opt{D}{T$_1$ is opened and the content of bucket is poured in the tub.}
\end{qoptions}
\end{bankq}
\begin{bankq}[id=SRGBT12-P1-CHEM-S1-Q34, source={SRGBT - Test 12 (10 Oct 2021)}, paper=1, qno=34, type=multi, subtopic={pH}, answer={A}, answer_src=key, difficulty=5, diff_parts={K2 S1 I1 C1 F1}, diff_note={concepts: standard problem in this chapter, ph | steps: apply the relevant results; compute}, group=SRGBT12-P1-CHEM-S1-G33-34]
A pH meter is immersed in the bucket of the original set up shown. It shows a reading of 7. Now S$_1$ is the first opened and all content is allowed to drain out into the bucket. pH of the solution in the bucket is measured again. Similarly S$_2$, S$_3$ and S$_4$ are sequentially opened and the pH of the solution in the bucket is measured each time after complete drainage. Opening which stop cock would cause maximum pH change (increase/decrease) compared to previous pH in bucket.
\begin{qoptions}
\opt{A}{S$_1$}
\opt{B}{S$_2$}
\opt{C}{S$_3$}
\opt{D}{S$_4$}
\end{qoptions}
\end{bankq}

\section*{Nurture SRG papers}
\begin{bankq}[id=NSRG02-P1-CHEM-S1-Q01, source={Nurture SRG - Test 02 (06 Dec 2023)}, paper=1, qno=1, type=single, subtopic={Hydrolysis}, answer={C}, answer_src=solved, working={Acetic acid has the lowest pH after the strong acid, ammonium acetate is near neutral, sodium acetate is basic and NaOH most basic: $V<II<III<I<IV$}, difficulty=3, diff_parts={K1 S2 I1 C0 F0}, diff_note={concepts: pH of different solutions | steps: pH of the salts; strong acid and base; rank}]
The pH value of 0.1 M solutions of $\mathrm{CH_3COONa}$ (I), $\mathrm{CH_3COOH}$ (II), $\mathrm{CH_3COONH_4}$ (III), NaOH (IV) and HCl (V) is in the order
\begin{qoptions}
\opt{A}{I $<$ II $<$ III $<$ IV $<$ V}
\opt{B}{V $<$ IV $<$ III $<$ II $<$ I}
\opt{C}{V $<$ II $<$ III $<$ I $<$ IV}
\opt{D}{V $<$ II $<$ I $<$ III $<$ IV}
\end{qoptions}
\end{bankq}
\begin{bankq}[id=NSRG02-P1-CHEM-S1-Q02, source={Nurture SRG - Test 02 (06 Dec 2023)}, paper=1, qno=2, type=single, subtopic={pH}, answer={C}, answer_src=solved, working={$[\mathrm{H^+}]=\sqrt{K_w}$, so $\mathrm{pH}=\tfrac12(14-\log2.5)=6.8$}, difficulty=4, diff_parts={K1 S2 I1 C1 F0}, diff_note={concepts: pH of a neutral solution at another temperature | steps: hydrogen ion concentration; logarithm; pH}]
What if the pH of a neutral solution at $37^\circ$C, where $K_w$ equals $2.5\times10^{-14}$ ? ($\log2.5=0.4$)
\begin{qoptions}
\opt{A}{7.0}
\opt{B}{13.6}
\opt{C}{6.8}
\opt{D}{6.6}
\end{qoptions}
\end{bankq}
\begin{bankq}[id=NSRG02-P1-CHEM-S1-Q03, source={Nurture SRG - Test 02 (06 Dec 2023)}, paper=1, qno=3, type=single, subtopic={pH}, answer={C}, answer_src=solved, working={$[\mathrm{H^+}]=10^{-12}$ M, so in 1 mL there are $10^{-15}\times6.02\times10^{23}=6.02\times10^8$ ions}, difficulty=4, diff_parts={K1 S2 I1 C1 F0}, diff_note={concepts: number of hydronium ions in a small volume | steps: concentration from the pH; moles; number of ions}]
The number of hydronium ions in 1 mL of an aqueous solution of pH 12.0 at $25^\circ$C is
\begin{qoptions}
\opt{A}{0.01}
\opt{B}{$10^{-12}$}
\opt{C}{$6.02\times10^8$}
\opt{D}{$6.02\times10^{11}$}
\end{qoptions}
\end{bankq}
\begin{bankq}[id=NSRG02-P1-CHEM-S2-Q05, source={Nurture SRG - Test 02 (06 Dec 2023)}, paper=1, qno=5, type=multi, subtopic={pH}, answer={A,C}, answer_src=solved, working={First dissociation dominates so $[\mathrm{H^+}]\approx\sqrt{K_{a_1}C}$ and $[\mathrm{H^+}]\approx[\mathrm{H_2PO_4^-}]$, hence $[\mathrm{HPO_4^{2-}}]\approx K_{a_2}$}, difficulty=5, diff_parts={K1 S2 I2 C0 F1}, diff_note={concepts: dissociation of phosphoric acid | steps: dominant first step; concentration of the ions; relation of Ka2; check}]
If $K_{a_1}$, $K_{a_2}$ and $K_{a_3}$ be the first, second and third dissociation constant of $\mathrm{H_3PO_4}$ and $K_{a_1}\gg K_{a_2}\gg K_{a_3}$ which is/are correct :
\begin{qoptions}
\opt{A}{$\left[\mathrm{H^+}\right]\simeq\sqrt{K_{a_1}\left[\mathrm{H_3PO_4}\right]}$}
\opt{B}{$\left[\mathrm{H^+}\right]\simeq\left[\mathrm{HPO_4^{2-}}\right]$}
\opt{C}{$K_{a_2}\simeq\left[\mathrm{HPO_4^{2-}}\right]$}
\opt{D}{$\left[\mathrm{HPO_4^{2-}}\right]=\left[\mathrm{PO_4^{3-}}\right]$}
\end{qoptions}
\end{bankq}
\begin{bankq}[id=NSRG02-P1-CHEM-S2-Q06, source={Nurture SRG - Test 02 (06 Dec 2023)}, paper=1, qno=6, type=multi, subtopic={Solubility product}, answer={A,C,D}, answer_src=solved, working={Option A: $s=\sqrt{K_{sp}}=10^{-5}$ M (true); C: $s=K_{sp}/2=5\times10^{-11}$ M; D: $K=K_{sp}K_f=10^{-2}$ gives $s/(2-2s)=0.1$, $s=0.166$ M; B is not supported (conversion to AgBr)}, difficulty=5, diff_parts={K2 S2 I0 C1 F1}, diff_note={concepts: solubility product, complex formation | steps: solubility in water; in KBr; in AgNO3; in NH3 using K = Ksp Kf and solve}]
Which is/are correct statement(s) about the solubility of AgCl(s) ?\par
Given : $K_{sp(AgCl)}=10^{-10}$; $K_f\left[\mathrm{Ag(NH_3)_2^+}\right]=10^8$.
\begin{qoptions}
\opt{A}{Solubility of AgCl in pure water is $10^{-5}$ mole/litre}
\opt{B}{Solubility of AgCl in 2M KBr is $10^{-5}$ gm/litre}
\opt{C}{Solubility of AgCl in 2 M $\mathrm{AgNO_3}$ is $5\times10^{-11}$ M}
\opt{D}{Solubility of AgCl in 2M $\mathrm{NH_3}$ is 0.166 M}
\end{qoptions}
\end{bankq}
\begin{bankq}[id=NSRG02-P1-CHEM-S2-Q07, source={Nurture SRG - Test 02 (06 Dec 2023)}, paper=1, qno=7, type=multi, subtopic={Hydrolysis}, answer={B,C,D}, answer_src=solved, working={$h=\sqrt{K_h/C}$ with $K_h=K_w/K_b$: it grows with dilution, with smaller $K_b$ and with temperature}, difficulty=3, diff_parts={K1 S1 I1 C0 F1}, diff_note={concepts: hydrolysis of salt of weak base | steps: expression for the degree of hydrolysis; effect of dilution, Kb, temperature}]
Degree of hydrolysis (h) for a salt of strong acid and weak base is :
\begin{qoptions}
\opt{A}{Independent of dilution}
\opt{B}{Increases with dilution}
\opt{C}{Increases with decrease in $K_b$}
\opt{D}{Increases with increase in temperature}
\end{qoptions}
\end{bankq}
\begin{bankq}[id=NSRG02-P1-CHEM-S2-Q08, source={Nurture SRG - Test 02 (06 Dec 2023)}, paper=1, qno=8, type=multi, subtopic={pH}, answer={A,C}, answer_src=solved, working={$K_{a_1}=\dfrac{(10^{-3})^2}{0.1}=10^{-5}$, $K_{a_2}=[\mathrm{A^{2-}}]=10^{-12}$ so $pK_{a_2}-pK_{a_1}=7$; $[\mathrm{OH^-}]=10^{-11}$}, difficulty=7, diff_parts={K1 S3 I2 C1 F1}, diff_note={concepts: diprotic weak acid from the pH | steps: first step dominates; hydroxide ion; Ka1; Ka2 from the anion; pKa difference}]
$\mathrm{H_2A}$ is a weak diprotic acid. If the pH of 0.1M $\mathrm{H_2A}$ solution is 3 and concentration of $\mathrm{A^{2-}}$ is $10^{-12}$ at $25^\circ$C.\par
Select correct statement(s)
\begin{qoptions}
\opt{A}{$\left[\mathrm{H^+}\right]_{\text{total}}\approx\left[\mathrm{H^+}\right]$ from first step of ionization of acid $\mathrm{H_2A}$}
\opt{B}{Concentration of $\mathrm{OH^-}$ in solution is $10^{-3}$ M}
\opt{C}{The value of $K_{a_1}$ is nearly $10^{-5}$}
\opt{D}{$pK_{a_2}-pK_{a_1}=9$}
\end{qoptions}
\end{bankq}
\begin{bankq}[id=NSRG02-P1-CHEM-S2-Q10, source={Nurture SRG - Test 02 (06 Dec 2023)}, paper=1, qno=10, type=multi, subtopic={pH}, answer={A,B,C}, answer_src=solved, working={$pK_{In}=8$ so the range is 7 to 9; $75\%$ acid to $75\%$ base changes pH by $2\log3=0.96$; the range lies within the strong acid strong base jump; the ratio 2 gives $\mathrm{pH}=7.7$}, difficulty=7, diff_parts={K1 S3 I2 C1 F1}, diff_note={concepts: acid base indicator ranges | steps: pH range of the indicator; ratio of forms; pH change; suitability; pH for given ratio}]
In an acidic indicator HIn has ionization constant is $10^{-8}$. The acid form of indicator is yellow and alkaline form is red. Which is/are correct statements ?\par
(Given : $\log2=0.3$, $\log3=0.48$)
\begin{qoptions}
\opt{A}{The pH range of indicator is 7 to 9}
\opt{B}{Change in pH is 0.96 when 75 \% yellow colour change to 75 \% red colour}
\opt{C}{This indicator is suitable for the titration of strong acid vs. strong base}
\opt{D}{pH of indicator is 8.3 when ratio of acid form to alkaline form is 2.}
\end{qoptions}
\end{bankq}
\begin{bankq}[id=NSRG02-P1-CHEM-S2-Q11, source={Nurture SRG - Test 02 (06 Dec 2023)}, paper=1, qno=11, type=multi, subtopic={Buffers}, answer={A,C}, answer_src=solved, working={$\mathrm{NH_4Cl}$ with half as much NaOH gives $\mathrm{NH_3}$ and $\mathrm{NH_4^+}$; sodium acetate with half the moles of HCl gives equal acid and salt}, difficulty=6, diff_parts={K1 S2 I2 C1 F1}, diff_note={concepts: buffer from reaction mixtures | steps: reaction of the two components; species left; buffer composition}]
A buffer solution can be prepared from a mixture of
\begin{qoptions}
\opt{A}{$\mathrm{NH_4Cl}$ and NaOH in 2:1 mole ratio}
\opt{B}{$\mathrm{CH_3COONa}$ and HCl in 1:1 mole ratio}
\opt{C}{$\mathrm{CH_3COONa}$ and HCl in 2:1 mole ratio}
\opt{D}{$\mathrm{CH_3COONa}$ and HCl in 1:2 mole ratio}
\end{qoptions}
\end{bankq}
\begin{bankq}[id=NSRG02-P1-CHEM-S2-Q12, source={Nurture SRG - Test 02 (06 Dec 2023)}, paper=1, qno=12, type=multi, subtopic={Hydrolysis}, answer={A,C,D}, answer_src=solved, working={$K_h=K_w/K_2=2\times10^{-4}$, $h=\sqrt{K_h/C}=0.02$, $[\mathrm{OH^-}]=0.01$ so $pOH=2$; $[\mathrm{H_2CO_3}]=\dfrac{K_w}{K_1}\cdot\dfrac{[\mathrm{HCO_3^-}]}{[\mathrm{OH^-}]}=2.5\times10^{-9}$}, difficulty=8, diff_parts={K1 S3 I2 C2 F1}, diff_note={concepts: hydrolysis of sodium carbonate | steps: hydrolysis constant; degree of hydrolysis; hydroxide ion; pH; second hydrolysis; carbonic acid}]
$\mathrm{H_2CO_3}$ ionizes as :\par
$\mathrm{H_2CO_3}\rightleftharpoons\mathrm{H^+}+\mathrm{HCO_3^-}$; $K_1=4.0\times10^{-6}$\par
$\mathrm{HCO_3^-}\rightleftharpoons\mathrm{CO_3^{2-}}$; $K_2=5.0\times10^{-11}$\par
The correct information(s) related with 0.5 M $-$ $\mathrm{Na_2CO_3}$ solution is/are
\begin{qoptions}
\opt{A}{The degree of hydrolysis of $\mathrm{Na_2CO_3}$ is 0.02}
\opt{B}{pH of solution is 2.0}
\opt{C}{pOH of solution is 2.0}
\opt{D}{$\left[\mathrm{H_2CO_3}\right]\approx2.5\times10^{-9}$ M}
\end{qoptions}
\end{bankq}
\begin{bankq}[id=NSRG02-P1-CHEM-S3-Q01, source={Nurture SRG - Test 02 (06 Dec 2023)}, paper=1, qno=1, type=numerical, subtopic={pH}, answer={7.5}, answer_src=solved, working={$[\mathrm{D^+}]=\sqrt{10^{-15}}=10^{-7.5}$ so $p^D=7.5$}, difficulty=4, diff_parts={K1 S1 I1 C1 F1}, diff_note={concepts: pD of heavy water | steps: ionic product; pD from the square root}]
The ionic product of heavy water, $\mathrm{D_2O}$, is $10\times10^{-16}$ at $7^\circ$C. The $p^D$ value of pure heavy water at $7^\circ$C is
\end{bankq}
\begin{bankq}[id=NSRG02-P1-CHEM-S3-Q02, source={Nurture SRG - Test 02 (06 Dec 2023)}, paper=1, qno=2, type=numerical, subtopic={pH}, answer={4}, answer_src=solved, working={$[\mathrm{HCOO^-}]=10^{-3}$ M and formic acid is $\dfrac{1150}{46}=25$ M, so $x=\dfrac{10^{-3}}{25}\times100=0.004$ and $1000x=4$}, difficulty=8, diff_parts={K1 S3 I2 C2 F1}, diff_note={concepts: self ionisation of formic acid | steps: ion concentration; concentration of formic acid; fraction ionised; percentage; scale}]
The self-ionization constant for pure formic acid, $K=\left[\mathrm{HCOOH_2^+}\right]\left[\mathrm{HCOO^-}\right]$ has been estimated as $10^{-6}\,\mathrm{M^2}$ and the density of formic acid is 1.15 g/cm$^3$ at room temperature. If `x $\%$' of formic acid molecules in pure formic acid is converted to formate ions, then the value of `1000x' is
\end{bankq}
\begin{bankq}[id=NSRG02-P1-CHEM-S3-Q03, source={Nurture SRG - Test 02 (06 Dec 2023)}, paper=1, qno=3, type=numerical, subtopic={Solubility product}, answer={8}, answer_src=solved, working={$K_{sp}=4(5\times10^{-6})^3=5\times10^{-16}$; at $s=5\times10^{-4}$, $[\mathrm{OH^-}]^2=10^{-12}$, so $pOH=6$ and $pH=8$}, difficulty=6, diff_parts={K1 S2 I2 C1 F1}, diff_note={concepts: pH of a buffer from the solubility | steps: solubility product from the first solubility; hydroxide for the second; pH}]
The solubility of $\mathrm{Pb(OH)_2}$ in water is $5\times10^{-6}$M. Calculate the pH of buffer solution in which solubility of $\mathrm{Pb(OH)_2}$ is $5\times10^{-4}$ mol litre$^{-1}$.
\end{bankq}
\begin{bankq}[id=NSRG02-P1-CHEM-S3-Q04, source={Nurture SRG - Test 02 (06 Dec 2023)}, paper=1, qno=4, type=numerical, subtopic={pH}, answer={4}, answer_src=solved, working={First equivalence at 20 mL: $\mathrm{pH}_1=\tfrac12(4.6+8)=6.3$; second at 40 mL with salt concentration $0.04$ M: $\mathrm{pH}_2=7+\tfrac12(8+\log0.04)=10.3$; difference 4}, difficulty=7, diff_parts={K1 S3 I2 C1 F1}, diff_note={concepts: pH difference between equivalence points | steps: first equivalence from the graph; pH from the pKa values; concentration at the second equivalence; salt pH; difference}]
10 mL of $\mathrm{H_2A}$ (weak diprotic acid) solution is titrated against 0.1 M NaOH. $p^H$ of the solution is plotted against volume of strong base added and following observation is made.\par
\bfig{100.2}{NSRG02_p1_chem_s3_q04_1.png}
If pH of the solution at first equivalence point is $pH_1$ and at second equivalence point is $pH_2$.\par
Calculate the value of $(pH_2-pH_1)$ at $25^\circ$C\par
Given: For $\mathrm{H_2A}$, $pK_{a_1}=4.6$ and $pK_{a_2}=8$, $\log25=1.4$
\end{bankq}
\begin{bankq}[id=NSRG02-P1-CHEM-S3-Q05, source={Nurture SRG - Test 02 (06 Dec 2023)}, paper=1, qno=5, type=numerical, subtopic={pH}, answer={4.89}, answer_src=solved, working={$p_{\mathrm{SO_2}}=10^{-5}$ atm, dissolved $1.28\times10^{-5}$ M; with $K_a=10^{-2}$, $x^2+10^{-2}x-1.28\times10^{-7}=0$ gives $x=1.28\times10^{-5}$ M, $\mathrm{pH}\approx4.89$}, difficulty=8, diff_parts={K1 S3 I2 C2 F1}, diff_note={concepts: pH of rain from atmospheric sulphur dioxide | steps: partial pressure from ppm; dissolved concentration; dissociation equilibrium; solve; pH}]
The average concentration of $\mathrm{SO_2}$ in the atmosphere over a city on a certain day is 10 ppm, when the average temperature is 298 K. Given that the solubility of $\mathrm{SO_2}$ in water at 298 K is 1.28 moles litre$^{-1}$ and the $pK_a$ of $\mathrm{H_2SO_3}$ is 2.0, estimate the pH of rain on that day.
\end{bankq}

\section*{Enthusiast Pre-Score papers}
\begin{bankq}[id=EPS02-P1-CHEM-S1-Q03, source={Enthusiast Pre-Score - Test 02 (11 Nov 2020)}, paper=1, qno=3, type=single, subtopic={Hydrolysis}, answer={B}, answer_src=key, difficulty=6, diff_parts={K2 S2 I2 C1 F0}, diff_note={concepts: pH of a salt from a hydrolysis graph, read pKa from the graph | steps: read pKa from the graph; neutralisation; buffer; pH}]
The given graph is for salt hydrolysis (monovalent type) between pH of solution vs. log c (where c = concentration of salt)\par
\bfig{97.7}{EPS02_p1_chem_s1_q03_1.png}
If 45 ml, 0.2 M of above salt is titrated with 60 ml, 0.1 M of NaOH, then pH of resulting solution will be :
\begin{qoptions}
\opt{A}{7.7}
\opt{B}{8.3}
\opt{C}{5.7}
\opt{D}{6.3}
\end{qoptions}
\end{bankq}
\begin{bankq}[id=EPS02-P1-CHEM-S3-Q06, source={Enthusiast Pre-Score - Test 02 (11 Nov 2020)}, paper=1, qno=6, type=numerical, subtopic={Solubility product}, answer={8.89}, answer_src=key, difficulty=7, diff_parts={K2 S2 I2 C1 F1}, diff_note={concepts: selective precipitation of hydroxides, solubility product | steps: Ksp of each hydroxide; hydroxide range; pH}]
If a solution contains 0.2 M each of Mg$^{+2}$ and Mn$^{+2}$ ions, to what pH must this solution be made so as to maximise precipitation of Mn(OH)$_2$ and almost no precipitation of Mg(OH)$_2$\par
Ksp [Mn(OH)$_2$] = 1.6 $\times$ 10$^{-13}$ ; Ksp [Mg(OH)$_2$] = 1.2 $\times$ 10$^{-11}$\par
Given (log 6 = 0.78 and log 8 = 0.9)
\end{bankq}
\begin{bankq}[id=EPS02-P2-CHEM-S1-Q05, source={Enthusiast Pre-Score - Test 02 (11 Nov 2020)}, paper=2, qno=5, type=multi, subtopic={pH}, answer={A,B,C}, answer_src=key, difficulty=7, diff_parts={K2 S2 I2 C1 F1}, diff_note={concepts: neutralisation constant of a weak acid, ph | steps: relation of K with Ka; pH; salt hydrolysis; mixed neutralisation}]
The equilibrium constant (K$_C$) for reaction of weak acid HA with strong base NaOH is 10$^9$ at 25$^\circ$C then -
\begin{qoptions}
\opt{A}{Ionisation constant of HA is 10$^{-5}$ at 25$^\circ$C}
\opt{B}{At 25$^\circ$C, pH of 0.01 M aq. solution of HA is 3.5}
\opt{C}{At 25$^\circ$C, pH of 0.1 M aq. solution of NaA is 9}
\opt{D}{If K$_b$ of weak acid BOH is 10$^{-4}$ at 25$^\circ$C, then equilibrium constant for neutralisation reaction between HA and BOH will be 10$^{-5}$ at 25$^\circ$C.}
\end{qoptions}
\end{bankq}
\begin{bankq}[id=EPS02-P2-CHEM-S3-Q05, source={Enthusiast Pre-Score - Test 02 (11 Nov 2020)}, paper=2, qno=5, type=numerical, subtopic={Solubility product}, answer={2.70}, answer_src=key, difficulty=6, diff_parts={K2 S2 I1 C1 F1}, diff_note={concepts: ion concentration from solubility product, solubility expression | steps: solubility expression; solve; value}]
A saturated solution of B(OH)$_3$ is prepared in pure water at 25$^\circ$C if Ksp of B(OH)$_3$ is 2.7 $\times$ 10$^{-33}$ then concentration of B$^{+3}$ is y $\times$ 10$^{-12}$ calculate value of y will be -\par
B(OH)$_3$(s) $\rightleftharpoons$ B$^{+3}$(aq.) + 3OH$^-$(aq.)
\end{bankq}

\section*{TOAS papers}
\begin{bankq}[id=TOAS12-P1-CHEM-S2-Q04, source={TOAS - Test 12 (10 Jul 2022)}, paper=1, qno=4, type=numerical, subtopic={pH}, answer={2.00}, answer_src=key, difficulty=2, diff_parts={K1 S1 I0 C0 F1}, diff_note={concepts: degree of dissociation of a weak base | steps: [OH-] from pH 11; alpha = [OH-]/C}]
Determine degree of dissociation (in \%) of 0.05 M NH$_3$ at 25$^\circ$C in a solution of pH = 11.
\end{bankq}
\begin{bankq}[id=TOAS14-P1-CHEM-S1-Q01, source={TOAS - Test 14 (07 Aug 2022)}, paper=1, qno=1, type=multi, subtopic={Buffers}, answer={A,C,D}, answer_src=key, difficulty=3, diff_parts={K1 S2 I0 C0 F1}, diff_note={concepts: Henderson equation for a phosphate buffer | steps: initial pH from pKa2 and the ratio 2:1; after HCl the ratio becomes 1:2; after NaOH the ratio is 5:1}]
A one litre solution contains 0.04 M Na$_2$HPO$_4$ and 0.02 M NaH$_2$PO$_4$. For H$_3$PO$_4$ : $\text{p}K_{a_1}=2.14$ ; $\text{p}K_{a_2}=6.84$ ; $\text{p}K_{a_3}=12.3$. The correct statement (s) is /are :\par [Assume volume of solution is constant]\par (use : $\log2=0.3$, $\log5=0.7$) :
\begin{qoptions}
\opt{A}{pH of solution is 7.14}
\opt{B}{On adding 0.02 moles of HCl to initial solution pH increases by 0.6 units}
\opt{C}{On adding 0.02 moles of HCl to initial solution pH decreases by 0.6 units}
\opt{D}{On adding 0.01 mol NaOH to initial solution pH of solution becomes 7.54}
\end{qoptions}
\end{bankq}
\begin{bankpassage}[id=TOAS16-P2-CHEM-S3-G3-4, source={TOAS - Test 16 (04 Sep 2022)}, questions={3-4}]
A weak base BOH is titrated with a strong acid HA. When 10 ml of HA is added, the pH is 10.3 and when 25 ml of HA is added the pH is 9.3.\par
In another experiment 40g of the salt BA is dissolved in water to make 1L aq. solution. The pH of solution is found to be 5.15. If salt BA forms NaCl type crystal having $r_{\text{B}^+}=0.8$\AA\ and $r_{\text{A}^-}=1.2$\AA\ ($\log2=0.3$, $N_A=6\times10^{23}$).
\end{bankpassage}
\begin{bankq}[id=TOAS16-P2-CHEM-S3-Q03, source={TOAS - Test 16 (04 Sep 2022)}, paper=2, qno=3, type=numerical, subtopic={Hydrolysis}, answer={4}, answer_src=solved, working={Buffer: pOH$=$pKb$+\log(\text{salt}/\text{base})$; with $n=30c$ the 10 mL point gives pOH$=3.7=$pKb$+\log\tfrac12$, so pKb$=4.0$}, difficulty=4, diff_parts={K1 S2 I1 C0 F1}, diff_note={concepts: titration of a weak base | steps: two buffer equations; mole relation; pKb}, group=TOAS16-P2-CHEM-S3-G3-4]
Find p$K_b$ of BOH :
\end{bankq}
\begin{bankq}[id=TOAS16-P2-CHEM-S3-Q04, source={TOAS - Test 16 (04 Sep 2022)}, paper=2, qno=4, type=numerical, subtopic={Hydrolysis}, answer={8.33}, answer_src=key, difficulty=6, diff_parts={K2 S2 I1 C1 F1}, diff_note={concepts: salt hydrolysis, rock-salt density | steps: pH gives concentration and molar mass; edge from the radii; density}, group=TOAS16-P2-CHEM-S3-G3-4]
Find density of BA crystal in g/ml.
\end{bankq}
\begin{bankq}[id=TOAS17-P1-CHEM-S1-Q07, source={TOAS - Test 17 (18 Sep 2022)}, paper=1, qno=7, type=multi, subtopic={pH}, answer={A,B,D}, answer_src=key, difficulty=3, diff_parts={K1 S2 I0 C0 F1}, diff_note={concepts: mixture of strong and weak acids | steps: pH from HCl; degree of dissociation of HA from the common ion effect; concentration of the anion}]
1L of 0.2 M HCl (aq.) is mixed with 1 L of 0.2 M weak monoprotic acid HA ($K_a=5\times10^{-6}$) at $25^\circ$C.\par
Select the correct options :
\begin{qoptions}
\opt{A}{pH of the final solution will be 1.}
\opt{B}{\% dissociation of HA in final solution is approximately $5\times10^{-3}$ \%}
\opt{C}{Conc. of A$^-$ in final solution will be $1\times10^{-3}$ M}
\opt{D}{Conc. of A$^-$ in final solution will be $5\times10^{-6}$ M}
\end{qoptions}
\end{bankq}
\begin{bankq}[id=TOAS17-P2-CHEM-S2-Q03, source={TOAS - Test 17 (18 Sep 2022)}, paper=2, qno=3, type=single, subtopic={Buffers}, answer={476}, answer_src=key, difficulty=2, diff_parts={K1 S1 I0 C1 F0}, diff_note={concepts: buffer pH at the half equivalence point | steps: pH = pKa; value times 100}]
50 mL of 0.1 M CH$_3$COOH is being titrated against 0.1 M NaOH. When 25 mL of NaOH has been added, the pH of the solution will be \underline{\hspace{1.2cm}} $\times10^{-2}$. (Nearest integer)
\par\noindent(Given : p$K_a$ (CH$_3$COOH) = 4.76)\par
\noindent$\log2=0.30$\par\noindent$\log3=0.48$\par\noindent$\log5=0.69$\par\noindent$\log7=0.84$\par\noindent$\log11=1.04$
\end{bankq}
\begin{bankq}[id=TOAS18-P1-CHEM-S3-Q05, source={TOAS - Test 18 (30 Oct 2022)}, paper=1, qno=5, type=numerical, subtopic={pH}, answer={8}, answer_src=solved, working={pOH $=4.5$: $K_b=[\mathrm{OH^-}]^2/C=10^{-9}/0.1=10^{-8}$, $\mathrm{p}K_b=8$}, difficulty=2, diff_parts={K1 S1 I0 C0 F1}, diff_note={concepts: pH of a weak base | steps: hydroxide concentration from the pH; Kb and pKb}]
0.1 M aqueous solution of a weak base BOH has pH = 9.5. Find p$K_b$ of the weak base, BOH ?
\end{bankq}

\section*{TAAS papers}
\begin{bankpassage}[id=TAAS02-P2-CHEM-S3-G3-4, source={TAAS - Test 02 (30 Jul 2023)}, questions={3-4}]
The solubility of any solid in water may be classified as

Molecular solubility : $\mathrm{AB(s)}\overset{\mathrm{aq.}}{\rightleftharpoons}\mathrm{AB(aq)}$

Ionic solubility : $\mathrm{AB(s)}\overset{\mathrm{aq.}}{\rightleftharpoons}\mathrm{A^+(aq)+B^-(aq)}$

For a substance AB(s) the following thermodynamics informations are available at 300 K :

Molecular solubility : $\Delta H^\circ=+20$ kcal/mol; $\Delta S^\circ=+40$ cal/K-mole

Ionic solubility : $\Delta H^\circ=-25$ kcal/mol; $\Delta S^\circ=-50$ cal/K-mole.
\end{bankpassage}
\begin{bankq}[id=TAAS02-P2-CHEM-S3-Q03, source={TAAS - Test 02 (30 Jul 2023)}, paper=2, qno=3, type=numerical, subtopic={Solubility product}, answer={3}, answer_src=key, difficulty=3, diff_parts={K1 S1 I1 C0 F1}, diff_note={concepts: spontaneity of dissolution routes | steps: delta G for each process; which occur}, group=TAAS02-P2-CHEM-S3-G3-4]
How many of the following statement(s) are \textbf{incorrect} for the solubility of AB(s) in water at 300 K :

(1) It undergoes molecular solubility

(2) It undergoes ionic solubility

(3) It undergoes ionic as well as molecular solubility

(4) It is completely insoluble in water.
\end{bankq}
\begin{bankq}[id=TAAS02-P2-CHEM-S3-Q04, source={TAAS - Test 02 (30 Jul 2023)}, paper=2, qno=4, type=numerical, subtopic={Solubility product}, answer={30}, answer_src=key, difficulty=4, diff_parts={K1 S2 I1 C0 F1}, diff_note={concepts: equilibrium constant from thermodynamic data | steps: combine the two processes; delta G; ln K}, group=TAAS02-P2-CHEM-S3-G3-4]
Calculate the value of $\ln K_{eq}$ at 300 K for the reaction

$\mathrm{AB(aq)\rightleftharpoons A^+(aq)+B^-(aq)}$
\end{bankq}

\section*{TMAS papers}
\begin{bankq}[id=TMAS04-P1-CHEM-S1-Q03, source={TMAS - Test 04 (19 Nov 2023)}, paper=1, qno=3, type=single, subtopic={Hydrolysis}, answer={A}, answer_src=key, difficulty=4, diff_parts={K1 S2 I1 C1 F0}, diff_note={concepts: neutralisation enthalpy and hydrolysis | steps: acid strength from enthalpy; conjugate base strength; pH}]
The enthalpy of neutralisation of four acids HA, HB, HC and HD with KOH are $-10.2$, $-11.4$, $-12.3$ and $-13.1$ kcal/mole respectively. Which of the following aqueous salt solutions has maximum pH ?
\begin{qoptions}
\opt{A}{$0.1$ M NaA}
\opt{B}{$0.1$ M NaB}
\opt{C}{$0.1$ M NaC}
\opt{D}{$0.1$ M NaD}
\end{qoptions}
\end{bankq}
\begin{bankq}[id=TMAS04-P1-CHEM-S3-Q03, source={TMAS - Test 04 (19 Nov 2023)}, paper=1, qno=3, type=numerical, subtopic={Solubility product}, answer={19.50}, answer_src=key, difficulty=7, diff_parts={K2 S2 I2 C1 F1}, diff_note={concepts: complex dissociation, solubility | steps: combine constants; equilibrium expression; solve concentration}]
Given :

$[\mathrm{Ag(NH_3)_2}^+](\mathrm{aq})\rightleftharpoons\mathrm{Ag}^+(\mathrm{aq})+2\mathrm{NH_3(aq)}$, $K_C=6\times10^{-7}$ and $K_{sp}$ of AgCl $=2.4\times10^{-10}$ at $298$ K

The concentration of the complex $[\mathrm{Ag(NH_3)_2}^+]$ in resultant solution, when solid AgCl is added to aqueous solution having $1$M ammonia to form saturated solution is $P\times10^{-3}$ M. Find value of $P$ ?
\end{bankq}
\begin{bankq}[id=TMAS04-P2-CHEM-S1-Q01, source={TMAS - Test 04 (19 Nov 2023)}, paper=2, qno=1, type=multi, subtopic={Buffers}, answer={B,C}, answer_src=key, difficulty=6, diff_parts={K1 S2 I2 C1 F1}, diff_note={concepts: buffer and pH calculations | steps: moles of each; excess acid or base; pH; check each}]
Select the correct statement(s) when $100$ ml of $0.1$ M $\mathrm{CH_3COOH}$ aqueous solution is mixed with following solutions at $25^\circ$C. [Assume volume of solid to be negligible]

[Given : $K_a\ \mathrm{CH_3COOH}=1\times10^{-5}$ M, $\log2=0.3$]
\begin{qoptions}
\opt{A}{On mixing with $50$ ml, $0.2$ M NaOH(aq), it will form a buffer solution having pH $=5$.}
\opt{B}{On mixing with $0.01$ mole of solid KOH, it will form an aqueous solution having pH $=9$}
\opt{C}{On mixing with $100$ ml, $0.1$ M HCl(aq), it will form an aqueous solution having pH $=1.3$}
\opt{D}{On mixing with $10$ ml, $0.2$ M NaOH(aq), it will form a buffer solution having pH $=4.6$.}
\end{qoptions}
\end{bankq}
\begin{bankq}[id=TMAS04-P2-CHEM-S3-Q02, source={TMAS - Test 04 (19 Nov 2023)}, paper=2, qno=2, type=numerical, subtopic={pH}, answer={8}, answer_src=key, difficulty=5, diff_parts={K1 S2 I1 C1 F1}, diff_note={concepts: isohydric solutions | steps: hydrogen ion concentration; weak acid; Ka}]
Isohydric solutions are the solutions having the same concentration of hydronium ion. If $0.2$ M HA (monobasic acid) solution is isohydric with $4\times10^{-4}$ M HCl solution, then $K_a$ for HA is $P\times10^{-7}$. Value of `P' is :
\end{bankq}

\section*{TNAS papers}
\begin{bankq}[id=TNAS01-P1-CHEM-S1-Q09, source={TNAS - Test 01 (05 Feb 2023)}, paper=1, qno=9, type=single, subtopic={pH}, answer={A}, answer_src=key, difficulty=3, diff_parts={K1 S1 I1 C1 F0}, diff_note={concepts: ionisation constants | steps: combine the constants; value}]
Consider the ionisation

$\mathrm{H^+(aq)+H_3N^+CH_2COO^-(aq)\rightleftharpoons}$

$\mathrm{H_3N^+CH_2COOH(aq)\rightleftharpoons H_2NCH_2COOH(aq)+H^+(aq)}$

Assuming that the following acid dissociation constant apply to the ionisations :

{\small$\mathrm{H_3N^+CH_2COOH(aq)\rightleftharpoons H^+(aq)+H_3N^+CH_2COO^-(aq)}$}; $K_a=4\times10^{-3}$

$\mathrm{H_3N^+CH_2COO^-(aq)\rightleftharpoons H_2NCH_2COOH(aq)+H^+(aq)}$; $K_a=1.6\times10^{-10}$

Estimate a value for the equilibrium constant for the process

$\mathrm{H_3N^+CH_2COO^-(aq)\rightleftharpoons H_2NCH_2COOH(aq)}$ is
\begin{qoptions}
\opt{A}{$4\times10^{-8}$}
\opt{B}{$0.25\times10^{7}$}
\opt{C}{$6.4\times10^{-13}$}
\opt{D}{$0.25\times10^{13}$}
\end{qoptions}
\end{bankq}
\begin{bankq}[id=TNAS07-P1-CHEM-S1-Q13, source={TNAS - Test 07 (20 Aug 2023)}, paper=1, qno=13, type=single, subtopic={pH}, answer={B}, answer_src=key, difficulty=4, diff_parts={K1 S2 I1 C1 F0}, diff_note={concepts: exact hydrogen ion concentration of a strong acid | steps: charge balance; water ionisation; select equation}]
Which of the following equation must be used for the exact calculation of $[\mathrm{H^+}]$ of an aqueous HCl solution at any concentration.
\begin{qoptions}
\opt{A}{$[\mathrm{H^+}]=[\mathrm{HCl}]_0+K_W$}
\opt{B}{$[\mathrm{H^+}]=[\mathrm{HCl}]_0+\dfrac{K_w}{[\mathrm{H^+}]}$}
\opt{C}{$[\mathrm{H^+}]=[\mathrm{HCl}]_0+\dfrac{K_w}{[\mathrm{OH^-}]}$}
\opt{D}{$[\mathrm{H^+}]=[\mathrm{HCl}]_0-\dfrac{K_w}{[\mathrm{OH^-}]}$}
\end{qoptions}
\end{bankq}
\begin{bankq}[id=TNAS07-P2-CHEM-S1-Q04, source={TNAS - Test 07 (20 Aug 2023)}, paper=2, qno=4, type=multi, subtopic={pH}, answer={A,B,D}, answer_src=key, difficulty=6, diff_parts={K1 S3 I1 C1 F1}, diff_note={concepts: weak acid from the pH | steps: hydrogen ion concentration; degree of dissociation; Ka; pOH; effect of dilution}]
pH of 0.1 M weak monoprotic acid (HA) is 1.7 at $25^\circ$C. ($\log2=0.30$, $\log10=1$)
\begin{qoptions}
\opt{A}{Acid dissociation constant of HA is $5\times10^{-3}$.}
\opt{B}{Degree of dissociation of 0.1 M HA solution is 20\%.}
\opt{C}{pOH of 0.1 M HA solution is 11.3.}
\opt{D}{Concentration of $\mathrm{H^+}$ ion in solution decreases on dilution of solution.}
\end{qoptions}
\end{bankq}

\section*{Star Review (Other) papers}
\begin{bankq}[id=STAR03-P1-CHEM-S4-Q07, source={Star Review (Other) - Test 03 (04 Dec 2023)}, paper=1, qno=7, type=numerical, subtopic={Solubility product}, answer={10}, answer_src=key, difficulty=5, diff_parts={K1 S2 I1 C1 F1}, diff_note={concepts: solubility of hydroxide in sodium hydroxide | steps: common ion; solubility product; hydroxide concentration; value of C}]
The molar solubility of $\mathrm{Zn(OH)_2}$ in $C\times10^{-2}$ M NaOH solution is $2\times10^{-18}$ M. The value of $C$ is \underline{\hspace{1cm}}(Nearest integer)\par
(Given : The solubility product of $\mathrm{Zn(OH)_2}$ is $2\times10^{-20}$)
\end{bankq}

\section*{Nurture (General) papers}
\begin{bankq}[id=NUR01-P1-CHEM-S2-Q01, source={Nurture (General) - Test 01 (12 Jun 2021)}, paper=1, qno=1, type=single, subtopic={Hydrolysis}, answer={D}, answer_src=solved, working={HCOOK 0.1 M: 2x0.1=0.2 particles; pH=7+(pKa+log c)/2=8.5; H2 formed at anode by formate oxidation}, difficulty=7, diff_parts={K3 S2 I2 C1 F0}, diff_note={concepts: match for potassium formate solution, hydrolysis, particle count | steps: particle count; electrolysis products; pH of salt; match}]
Which of the following option is correctly matched for 0.1 M HCOOK(aq.)
\begin{qoptions}
\opt{A}{(I),(B),(S)}
\opt{B}{(IV), (B), (S)}
\opt{C}{(III),(D), (Q)}
\opt{D}{(III),(C),(P)}
\end{qoptions}
\end{bankq}
\begin{bankq}[id=NUR01-P1-CHEM-S2-Q03, source={Nurture (General) - Test 01 (12 Jun 2021)}, paper=1, qno=3, type=single, subtopic={Buffers}, answer={A}, answer_src=solved, working={Acidic buffer pH=5+log(0.1/0.2)=4.7; particles 0.4; H2 at cathode}, difficulty=7, diff_parts={K3 S2 I2 C1 F0}, diff_note={concepts: match for an acetate buffer, buffers, particle count | steps: particle count; buffer pH; electrolysis; match}]
Select the correct match for aq. solution having 0.1 M $\mathrm{CH_3COONa}$ and 0.2 M $\mathrm{CH_3COOH}$.
\begin{qoptions}
\opt{A}{(II),(B),(Q)}
\opt{B}{(II),(A),(Q)}
\opt{C}{(IV),(A),(S)}
\opt{D}{(I),(C),(R)}
\end{qoptions}
\end{bankq}
\begin{bankq}[id=NUR04-P1-CHEM-S2-Q10, source={Nurture (General) - Test 04 (11 Feb 2024)}, paper=1, qno=10, type=multi, subtopic={Buffers}, answer={C,D}, answer_src=solved, working={HNO3 with CH3COONa gives CH3COOH+CH3COONa when the acid is limiting; the acid and its salt directly}, difficulty=4, diff_parts={K1 S2 I1 C0 F1}, diff_note={concepts: buffer from mixtures | steps: reactions on mixing; weak acid and conjugate base; check pairs}]
Aqueous solutions of HNO$_3$, KOH, CH$_3$COOH, and CH$_3$COONa of identical concentrations are provided. The pair(s) of solutions which form a buffer upon mixing is (are)
\begin{qoptions}
\opt{A}{HNO$_3$ and CH$_3$COOH}
\opt{B}{KOH and CH$_3$COONa}
\opt{C}{HNO$_3$ and CH$_3$COONa}
\opt{D}{CH$_3$COOH and CH$_3$COONa}
\end{qoptions}
\end{bankq}
\begin{bankq}[id=NUR04-P1-CHEM-S3-Q13, source={Nurture (General) - Test 04 (11 Feb 2024)}, paper=1, qno=13, type=numerical, subtopic={Hydrolysis}, answer={3}, answer_src=solved, working={Basic solutions: KCN, K2CO3 and LiCN}, difficulty=4, diff_parts={K1 S2 I1 C0 F1}, diff_note={concepts: solutions turning litmus blue | steps: hydrolysis of each salt; select the basic ones; count}]
Amongst the following, the total number of compounds whose aqueous solution turns red litmus paper blue is\par
KCN \quad K$_2$SO$_4$ \quad (NH$_4$)$_2$C$_2$O$_4$ \quad NaCl \quad Zn(NO$_3$)$_2$ \quad FeCl$_3$ \quad K$_2$CO$_3$ \quad NH$_4$NO$_3$ \quad LiCN
\end{bankq}
\end{document}
